% Complete reviewed Unit 6; editorial master: chapters/ch06/answer-key.md
% Markdown SHA-256: f32c7b7a4f9ca4d83c1694a268141ac4d49b3f5f9df20b98967c6fd1918e16a0
% One-use initial assembly followed by complete content and page review.
\newcommand{\runningtitle}{Unit 6. Answer key}
% Layout inherited from the reviewed previous course; set a running title per unit.
% This is a small hand-maintained TeX layout, not a conversion pipeline.
\RequirePackage{fix-cm}
\documentclass[11pt,twoside]{article}
\usepackage{fontspec}
\setmainfont{Cambria}[
  Renderer=HarfBuzz,
  BoldFont=cambriab.ttf,
  ItalicFont=cambriai.ttf,
  BoldItalicFont=cambriaz.ttf
]
\setmonofont{lmmono10-regular.otf}[BoldFont=lmmonolt10-bold.otf]
\usepackage[paperwidth=165mm,paperheight=241mm,inner=21mm,outer=19mm,
  top=22mm,bottom=20mm,headheight=14pt,headsep=7mm,footskip=10mm]{geometry}
\usepackage{xcolor,array,booktabs,longtable,tabularx}
\usepackage{enumitem,fancyhdr}
\usepackage[unicode,hidelinks]{hyperref}
\definecolor{ink}{HTML}{000000}
\definecolor{muted}{HTML}{4C5558}
\setlength{\parindent}{0pt}
\setlength{\parskip}{3.5pt}
\AtBeginDocument{\fontsize{11}{14.3}\selectfont}
\setlength{\tabcolsep}{4pt}
\renewcommand{\arraystretch}{1.10}
\setlist[itemize]{leftmargin=1.3em,label=\textbullet,itemsep=2pt,topsep=3pt,parsep=0pt}
\setlist[enumerate]{leftmargin=1.8em,itemsep=3pt,topsep=4pt,parsep=0pt}
\widowpenalty=10000
\clubpenalty=10000
\emergencystretch=1em
\raggedbottom
\pagestyle{fancy}
\fancyhf{}
\providecommand{\runningtitle}{Old English}
\fancyhead[LE,RO]{\small\itshape \runningtitle}
\fancyhead[LO,RE]{\small Old English}
\fancyfoot[LE,RO]{\small\thepage}
\renewcommand{\headrulewidth}{0.3pt}
\newcommand{\booktitle}[2]{
  {\fontsize{9}{11}\selectfont\color{muted}#1\par}
  \vspace{3pt}
  {\fontsize{20}{24}\selectfont\color{ink}#2\par}
  \vspace{6pt}
}
\newcommand{\source}[1]{\par{\footnotesize\color{muted}#1\par}}
\newcommand{\readingtext}[1]{
  \par\addvspace{5pt}
  \begingroup\leftskip=1.5em
    \fontsize{12}{16}\selectfont #1\par
  \endgroup\addvspace{5pt}
}
\makeatletter
\renewcommand{\section}{\@startsection{section}{1}{\z@}
  {-2.4ex \@plus -0.6ex \@minus -0.2ex}{0.8ex \@plus 0.2ex}
  {\normalfont\fontsize{12}{15}\selectfont\bfseries\color{ink}}}
\renewcommand{\subsection}{\@startsection{subsection}{2}{\z@}
  {-1.8ex \@plus -0.5ex \@minus -0.2ex}{0.5ex \@plus 0.2ex}
  {\normalfont\fontsize{11}{14}\selectfont\bfseries\color{ink}}}
\makeatother

\hyphenpenalty=10000
\exhyphenpenalty=10000
\emergencystretch=3em
\newcommand{\itempara}[2]{\par\addvspace{3pt}\noindent\hangindent=1.8em\hangafter=1\makebox[1.8em][l]{#1.}#2\par}
\newcommand{\reserve}[1]{\par\penalty-100\begingroup\dimen0=\pagegoal\advance\dimen0 by-\pagetotal\ifdim\dimen0<#1\newpage\fi\endgroup}
\begin{document}
\addfontfeatures{Ligatures=TeXOff}
\raggedright
\booktitle{UNIT 6}{Answer key}

% mdparagraph:P001:begin
The Old English sentences below are models. Accept another defensible word order or a stated alternative when it preserves the requested sense, construction, agreement and reference. A valid form can nevertheless miss an expressly selected target: distinguish that result from impossible Old English. In particular, documented \textbf{-enne} infinitives and finite variants need not be linguistic errors, although a task may request the selected table forms; a bare purpose infinitive does not meet an explicit request for the inflected construction. [1,4]
% mdparagraph:P001:end

% mdparagraph:P002:begin
In the command tables, singular means one addressee and plural means two or more. In infinitive constructions, the finite verb carries person and number; the infinitive retains the case government of its verb. A form exercise asks for forms, without introducing every possible use of the verb. Preserve the printed vowel quantities and stem alternations.
% mdparagraph:P002:end

% mdparagraph:P003:begin
For translations, judge the relations and meaning, not a unique English wording. A wish or purpose does not by itself assert fulfillment. Keep the reading’s supplied phrases available when answering its questions; their appearance here does not require independent production of their syntax.
% mdparagraph:P003:end

\reserve{7\baselineskip}
% mdparagraph:P004:begin
\section*{A. Indicative systems}
% mdparagraph:P004:end

\par\begin{minipage}{\linewidth}\raggedright
% item: NEW-U06-A01
% key: KEY-NEW-U06-A01
% mdparagraph:P005:begin
\textbf{A01.}
% mdparagraph:P005:end

% mdparagraph:P006:begin
\textbf{Target: the complete indicative system of witan.}
% mdparagraph:P006:end

% mdtable:T01:begin
% course-table-style: 2026-10
% course-panel:T01-P01:begin
\par\begingroup
\fontsize{10.5}{13.2}\selectfont
\setlength{\LTleft}{0pt}\setlength{\LTright}{\fill}
\noindent
\begin{tabular}{@{}>{\raggedright\arraybackslash}p{39.459mm}>{\raggedright\arraybackslash}p{39.459mm}>{\raggedright\arraybackslash}p{39.459mm}@{}}
\toprule
\textbf{Person and number} & \textbf{Present indicative} & \textbf{Preterite indicative} \\
\midrule
1 sg. & wāt & wiste \\
2 sg. & wāst & wistest \\
3 sg. & wāt & wiste \\
1–3 pl. & witon & wiston \\
\bottomrule
\end{tabular}
\par\endgroup
% course-panel:T01-P01:end
% mdtable:T01:end
% mdparagraph:P007:begin
Present knowledge, “know”; past knowledge, “knew.” \textbf{Wāt} serves both first and third singular, but \textbf{wāst} is second singular. The present plural is \textbf{witon}; the dental past is \textbf{wiste}, with \textbf{wistest} in the second singular and \textbf{wiston} in every plural person. A strong-looking present stem does not give the word past meaning. [1]
% mdparagraph:P007:end

\end{minipage}\par\addvspace{4pt}

\par\begin{minipage}{\linewidth}\raggedright
% item: NEW-U06-A02
% key: KEY-NEW-U06-A02
% mdparagraph:P008:begin
\textbf{A02.}
% mdparagraph:P008:end

% mdparagraph:P009:begin
\textbf{Target: the complete indicative system of cunnan.}
% mdparagraph:P009:end

% mdtable:T02:begin
% course-table-style: 2026-10
% course-panel:T02-P01:begin
\par\begingroup
\fontsize{10.5}{13.2}\selectfont
\setlength{\LTleft}{0pt}\setlength{\LTright}{\fill}
\noindent
\begin{tabular}{@{}>{\raggedright\arraybackslash}p{39.459mm}>{\raggedright\arraybackslash}p{39.459mm}>{\raggedright\arraybackslash}p{39.459mm}@{}}
\toprule
\textbf{Person and number} & \textbf{Present indicative} & \textbf{Preterite indicative} \\
\midrule
1 sg. & cann & cūðe \\
2 sg. & canst & cūðest \\
3 sg. & cann & cūðe \\
1–3 pl. & cunnon & cūðon \\
\bottomrule
\end{tabular}
\par\endgroup
% course-panel:T02-P01:end
% mdtable:T02:end
% mdparagraph:P010:begin
In the taught skill construction, the present means “know how to,” the past “knew how to.” Distinguish singular \textbf{cann/\allowbreak{}canst} from plural \textbf{cunnon}. The past has long \textbf{ū} and a dental consonant: \textbf{cūðe}, \textbf{cūðest}, \textbf{cūðon}. Do not use the present singular as a past simply because it resembles a strong preterite. [1]
% mdparagraph:P010:end

\end{minipage}\par\addvspace{4pt}

\par\begin{minipage}{\linewidth}\raggedright
% item: NEW-U06-A03
% key: KEY-NEW-U06-A03
% mdparagraph:P011:begin
\textbf{A03.}
% mdparagraph:P011:end

% mdparagraph:P012:begin
\textbf{Target: the complete indicative system of ðurfan.}
% mdparagraph:P012:end

% mdtable:T03:begin
% course-table-style: 2026-10
% course-panel:T03-P01:begin
\par\begingroup
\fontsize{10.5}{13.2}\selectfont
\setlength{\LTleft}{0pt}\setlength{\LTright}{\fill}
\noindent
\begin{tabular}{@{}>{\raggedright\arraybackslash}p{39.459mm}>{\raggedright\arraybackslash}p{39.459mm}>{\raggedright\arraybackslash}p{39.459mm}@{}}
\toprule
\textbf{Person and number} & \textbf{Present indicative} & \textbf{Preterite indicative} \\
\midrule
1 sg. & ðearf & ðorfte \\
2 sg. & ðearft & ðorftest \\
3 sg. & ðearf & ðorfte \\
1–3 pl. & ðurfon & ðorfton \\
\bottomrule
\end{tabular}
\par\endgroup
% course-panel:T03-P01:end
% mdtable:T03:end
% mdparagraph:P013:begin
The selected sense is “need to,” with “needed to” in the past. The singular present has \textbf{ðearf-}, the plural \textbf{ðurf-}, and the past \textbf{ðorft-}. Preserve the final \textbf{-t} of second-singular \textbf{ðearft} and the full past ending in \textbf{ðorftest}; the present and past plural stems are different. [1]
% mdparagraph:P013:end

\end{minipage}\par\addvspace{4pt}

\par\begin{minipage}{\linewidth}\raggedright
% item: NEW-U06-A04
% key: KEY-NEW-U06-A04
% mdparagraph:P014:begin
\textbf{A04.}
% mdparagraph:P014:end

% mdparagraph:P015:begin
\textbf{Target: the complete indicative system of durran.}
% mdparagraph:P015:end

% mdtable:T04:begin
% course-table-style: 2026-10
% course-panel:T04-P01:begin
\par\begingroup
\fontsize{10.5}{13.2}\selectfont
\setlength{\LTleft}{0pt}\setlength{\LTright}{\fill}
\noindent
\begin{tabular}{@{}>{\raggedright\arraybackslash}p{39.459mm}>{\raggedright\arraybackslash}p{39.459mm}>{\raggedright\arraybackslash}p{39.459mm}@{}}
\toprule
\textbf{Person and number} & \textbf{Present indicative} & \textbf{Preterite indicative} \\
\midrule
1 sg. & dearr & dorste \\
2 sg. & dearst & dorstest \\
3 sg. & dearr & dorste \\
1–3 pl. & durron & dorston \\
\bottomrule
\end{tabular}
\par\endgroup
% course-panel:T04-P01:end
% mdtable:T04:end
% mdparagraph:P016:begin
The selected meanings are “dare to” and “dared to.” The first and third singular use \textbf{dearr}, the second \textbf{dearst}, and the plural \textbf{durron}. The past changes to \textbf{dorst-}. A plural subject requires the plural form even when English “dare” does not visibly distinguish the persons. [1]
% mdparagraph:P016:end

\end{minipage}\par\addvspace{4pt}

\par\begin{minipage}{\linewidth}\raggedright
% item: NEW-U06-A05
% key: KEY-NEW-U06-A05
% mdparagraph:P017:begin
\textbf{A05.}
% mdparagraph:P017:end

% mdparagraph:P018:begin
\textbf{Target: the complete indicative system of sculan.}
% mdparagraph:P018:end

% mdtable:T05:begin
% course-table-style: 2026-10
% course-panel:T05-P01:begin
\par\begingroup
\fontsize{10.5}{13.2}\selectfont
\setlength{\LTleft}{0pt}\setlength{\LTright}{\fill}
\noindent
\begin{tabular}{@{}>{\raggedright\arraybackslash}p{39.459mm}>{\raggedright\arraybackslash}p{39.459mm}>{\raggedright\arraybackslash}p{39.459mm}@{}}
\toprule
\textbf{Person and number} & \textbf{Present indicative} & \textbf{Preterite indicative} \\
\midrule
1 sg. & sceal & sceolde \\
2 sg. & scealt & sceoldest \\
3 sg. & sceal & sceolde \\
1–3 pl. & sculon & sceoldon \\
\bottomrule
\end{tabular}
\par\endgroup
% course-panel:T05-P01:end
% mdtable:T05:end
% mdparagraph:P019:begin
In the assigned obligation sense, the present means “must,” the preterite “had to.” Contrast \textbf{sceal/\allowbreak{}scealt}, \textbf{sculon}, and past \textbf{sceold-}. The table retains \textbf{eo} in the past. Cook also records the shorter spelling represented by \textbf{scoldon}: it is a valid variant, but it differs from the selected table. The past need not mean hypothetical English “should.” [1]
% mdparagraph:P019:end

\end{minipage}\par\addvspace{4pt}

\par\begin{minipage}{\linewidth}\raggedright
% item: NEW-U06-A06
% key: KEY-NEW-U06-A06
% mdparagraph:P020:begin
\textbf{A06.}
% mdparagraph:P020:end

% mdparagraph:P021:begin
\textbf{Target: the complete indicative system of mugan.}
% mdparagraph:P021:end

% mdtable:T06:begin
% course-table-style: 2026-10
% course-panel:T06-P01:begin
\par\begingroup
\fontsize{10.5}{13.2}\selectfont
\setlength{\LTleft}{0pt}\setlength{\LTright}{\fill}
\noindent
\begin{tabular}{@{}>{\raggedright\arraybackslash}p{39.459mm}>{\raggedright\arraybackslash}p{39.459mm}>{\raggedright\arraybackslash}p{39.459mm}@{}}
\toprule
\textbf{Person and number} & \textbf{Present indicative} & \textbf{Preterite indicative} \\
\midrule
1 sg. & mæg & meahte \\
2 sg. & meaht & meahtest \\
3 sg. & mæg & meahte \\
1–3 pl. & magon & meahton \\
\bottomrule
\end{tabular}
\par\endgroup
% course-panel:T06-P01:end
% mdtable:T06:end
% mdparagraph:P022:begin
The selected present expresses ability; the preterite expresses past ability. The second-singular present \textbf{meaht} differs from first/\allowbreak{}third \textbf{mæg} and plural \textbf{magon}. The chosen past has \textbf{meaht-} throughout. \textbf{Mugan} and \textbf{magan} are citation labels for this system; the label is not an assigned usable infinitive. [1]
% mdparagraph:P022:end

\end{minipage}\par\addvspace{4pt}

\par\begin{minipage}{\linewidth}\raggedright
% item: NEW-U06-A07
% key: KEY-NEW-U06-A07
% mdparagraph:P023:begin
\textbf{A07.}
% mdparagraph:P023:end

% mdparagraph:P024:begin
\textbf{Target: the complete indicative system of mōtan.}
% mdparagraph:P024:end

% mdtable:T07:begin
% course-table-style: 2026-10
% course-panel:T07-P01:begin
\par\begingroup
\fontsize{10.5}{13.2}\selectfont
\setlength{\LTleft}{0pt}\setlength{\LTright}{\fill}
\noindent
\begin{tabular}{@{}>{\raggedright\arraybackslash}p{39.459mm}>{\raggedright\arraybackslash}p{39.459mm}>{\raggedright\arraybackslash}p{39.459mm}@{}}
\toprule
\textbf{Person and number} & \textbf{Present indicative} & \textbf{Preterite indicative} \\
\midrule
1 sg. & mōt & mōste \\
2 sg. & mōst & mōstest \\
3 sg. & mōt & mōste \\
1–3 pl. & mōton & mōston \\
\bottomrule
\end{tabular}
\par\endgroup
% course-panel:T07-P01:end
% mdtable:T07:end
% mdparagraph:P025:begin
Here the present expresses permission, “may/\allowbreak{}is allowed to,” and the past “was allowed to.” Keep long \textbf{ō} throughout. The second-singular present is \textbf{mōst}, but its past counterpart is \textbf{mōstest}. The dictionary also records other senses; substituting obligation would miss the permission target used in this unit. [1]
% mdparagraph:P025:end

\end{minipage}\par\addvspace{4pt}

\par\begin{minipage}{\linewidth}\raggedright
% item: NEW-U06-A08
% key: KEY-NEW-U06-A08
% mdparagraph:P026:begin
\textbf{A08.}
% mdparagraph:P026:end

% mdparagraph:P027:begin
\textbf{Target: the complete indicative system of dōn.}
% mdparagraph:P027:end

% mdtable:T08:begin
% course-table-style: 2026-10
% course-panel:T08-P01:begin
\par\begingroup
\fontsize{10.5}{13.2}\selectfont
\setlength{\LTleft}{0pt}\setlength{\LTright}{\fill}
\noindent
\begin{tabular}{@{}>{\raggedright\arraybackslash}p{39.459mm}>{\raggedright\arraybackslash}p{39.459mm}>{\raggedright\arraybackslash}p{39.459mm}@{}}
\toprule
\textbf{Person and number} & \textbf{Present indicative} & \textbf{Preterite indicative} \\
\midrule
1 sg. & dō & dyde \\
2 sg. & dēst & dydest \\
3 sg. & dēð & dyde \\
1–3 pl. & dōð & dydon \\
\bottomrule
\end{tabular}
\par\endgroup
% course-panel:T08-P01:end
% mdtable:T08:end
% mdparagraph:P028:begin
The ordinary meanings are “do” and “did.” Present \textbf{dō} and \textbf{dōð} have \textbf{ō}, but second/\allowbreak{}third singular have \textbf{ē}. The past is \textbf{dyde}, not a past made by adding an ending to the present stem. Its \textbf{y} is short. \textbf{Dōð} is the present form for all three plural persons, not an infinitive. [1]
% mdparagraph:P028:end

\end{minipage}\par\addvspace{4pt}

\par\begin{minipage}{\linewidth}\raggedright
% item: NEW-U06-A09
% key: KEY-NEW-U06-A09
% mdparagraph:P029:begin
\textbf{A09.}
% mdparagraph:P029:end

% mdparagraph:P030:begin
\textbf{Target: the complete indicative system of gān.}
% mdparagraph:P030:end

% mdtable:T09:begin
% course-table-style: 2026-10
% course-panel:T09-P01:begin
\par\begingroup
\fontsize{10.5}{13.2}\selectfont
\setlength{\LTleft}{0pt}\setlength{\LTright}{\fill}
\noindent
\begin{tabular}{@{}>{\raggedright\arraybackslash}p{39.459mm}>{\raggedright\arraybackslash}p{39.459mm}>{\raggedright\arraybackslash}p{39.459mm}@{}}
\toprule
\textbf{Person and number} & \textbf{Present indicative} & \textbf{Preterite indicative} \\
\midrule
1 sg. & gā & ēode \\
2 sg. & gǣst & ēodest \\
3 sg. & gǣð & ēode \\
1–3 pl. & gāð & ēodon \\
\bottomrule
\end{tabular}
\par\endgroup
% course-panel:T09-P01:end
% mdtable:T09:end
% mdparagraph:P031:begin
The ordinary meanings are “go” and “went.” Distinguish present \textbf{gā/\allowbreak{}gāð} from second/\allowbreak{}third singular \textbf{gǣst/\allowbreak{}gǣð}. The past uses the separate \textbf{ēod-} stem: \textbf{ēode}, \textbf{ēodest}, \textbf{ēodon}. Do not derive its past mechanically from \textbf{gān} or replace long \textbf{ǣ} in the selected present singulars with \textbf{ā}. [1]
% mdparagraph:P031:end

\end{minipage}\par\addvspace{4pt}

\par\begin{minipage}{\linewidth}\raggedright
% item: NEW-U06-A10
% key: KEY-NEW-U06-A10
% mdparagraph:P032:begin
\textbf{A10.}
% mdparagraph:P032:end

% mdparagraph:P033:begin
\textbf{Target: the complete indicative system of willan.}
% mdparagraph:P033:end

% mdtable:T10:begin
% course-table-style: 2026-10
% course-panel:T10-P01:begin
\par\begingroup
\fontsize{10.5}{13.2}\selectfont
\setlength{\LTleft}{0pt}\setlength{\LTright}{\fill}
\noindent
\begin{tabular}{@{}>{\raggedright\arraybackslash}p{39.459mm}>{\raggedright\arraybackslash}p{39.459mm}>{\raggedright\arraybackslash}p{39.459mm}@{}}
\toprule
\textbf{Person and number} & \textbf{Present indicative} & \textbf{Preterite indicative} \\
\midrule
1 sg. & wille & wolde \\
2 sg. & wilt & woldest \\
3 sg. & wille & wolde \\
1–3 pl. & willað & woldon \\
\bottomrule
\end{tabular}
\par\endgroup
% course-panel:T10-P01:end
% mdtable:T10:end
% mdparagraph:P034:begin
The selected meanings are “want/\allowbreak{}intend to” and “wanted/\allowbreak{}intended to.” Both first- and third-singular present are \textbf{wille}; second singular is \textbf{wilt}. The plural uses \textbf{willað}, while the past uses \textbf{wold-}. Do not add an ordinary third-singular dental ending to \textbf{wille}, or assume that \textbf{wolde} necessarily expresses a hypothetical “would.” [1]
% mdparagraph:P034:end

\end{minipage}\par\addvspace{4pt}

\par\begin{minipage}{\linewidth}\raggedright
% item: NEW-U06-A11
% key: KEY-NEW-U06-A11
% mdparagraph:P035:begin
\textbf{A11.}
% mdparagraph:P035:end

% mdparagraph:P036:begin
\textbf{Target: both complete present indicative series.}
% mdparagraph:P036:end

% mdtable:T11:begin
% course-table-style: 2026-10
% course-panel:T11-P01:begin
\par\begingroup
\fontsize{10.5}{13.2}\selectfont
\setlength{\LTleft}{0pt}\setlength{\LTright}{\fill}
\noindent
\begin{tabular}{@{}>{\raggedright\arraybackslash}p{39.459mm}>{\raggedright\arraybackslash}p{39.459mm}>{\raggedright\arraybackslash}p{39.459mm}@{}}
\toprule
\textbf{Person and number} & \textbf{bēo series} & \textbf{Earlier eom series} \\
\midrule
1 sg. & bēo & eom \\
2 sg. & bist & eart \\
3 sg. & bið & is \\
1–3 pl. & bēoð & sindon \\
\bottomrule
\end{tabular}
\par\endgroup
% course-panel:T11-P01:end
% mdtable:T11:end
% mdparagraph:P037:begin
There is no separate invented past paradigm for the bēo series: the earlier \textbf{wæs/\allowbreak{}wǣre/\allowbreak{}wǣron} system supplies the past. In the explicitly later situation of M12, \textbf{bið} means “will be”; the series can also express habitual or general truths. Its form alone does not require future meaning. These indicative cells are distinct from the earlier supplied wish forms in the creation reading. [1–2]
% mdparagraph:P037:end

\end{minipage}\par\addvspace{4pt}

\par\begin{minipage}{\linewidth}\raggedright
% item: NEW-U06-A12
% key: KEY-NEW-U06-A12
% mdparagraph:P038:begin
\textbf{A12.}
% mdparagraph:P038:end

% mdparagraph:P039:begin
\textbf{Target: tense and person within the indicative.}
% mdparagraph:P039:end

% mdtable:T12:begin
\par\begingroup
\fontsize{10.5}{13.2}\selectfont
\setlength{\LTleft}{0pt}\setlength{\LTright}{\fill}
\noindent
\begin{tabular}{@{}>{\raggedright\arraybackslash}p{26.48mm}>{\raggedright\arraybackslash}p{54.95mm}>{\raggedright\arraybackslash}p{36.94mm}@{}}
\toprule
\textbf{Form} & \textbf{With a third-singular subject} & \textbf{Also first singular?} \\
\midrule
wāt & Present: he knows & Yes \\
wiste & Preterite: he knew & Yes \\
cann & Present: he knows how/\allowbreak{}can & Yes \\
cūðe & Preterite: he knew how/\allowbreak{}could & Yes \\
gǣð & Present: he goes & No \\
ēode & Preterite: he went & Yes \\
dēð & Present: he does & No \\
dyde & Preterite: he did & Yes \\
\bottomrule
\end{tabular}
\par\endgroup

% mdtable:T12:end
% mdparagraph:P040:begin
\textbf{Wāt} and \textbf{cann} belong to preterite-present systems: their present forms resemble strong past forms historically, but their meaning here is present. Their actual pasts are \textbf{wiste} and \textbf{cūðe}. Shared spelling permits more than one person; an expressed subject selects the person in a sentence. The stated past-ability sense of \textbf{cūðe} need not be hypothetical. [1]
% mdparagraph:P040:end

\end{minipage}\par\addvspace{4pt}

\par\begin{minipage}{\linewidth}\raggedright
% item: NEW-U06-A13
% key: KEY-NEW-U06-A13
% mdparagraph:P041:begin
\textbf{A13.}
% mdparagraph:P041:end

% mdparagraph:P042:begin
\textbf{Target: distinguish missing forms from shared forms.} A dictionary needs a label under which to list a verb even when that label is not an attested infinitive. It therefore does not license an infinitive clause, an inflected infinitive, or a regular command pair. This applies to the reference labels for \textbf{durran, sculan, mugan/\allowbreak{}magan} and \textbf{mōtan}.
% mdparagraph:P042:end

% mdparagraph:P043:begin
A missing cell lacks an established usable form; syncretism means that two established cells have the same form. For example, \textbf{gāð} can be indicative or imperative in the appropriate construction. Some preterite-presents have documented command usages or substitutes, but these do not supply an ordinary complete imperative system for this exercise. Treating every dictionary label as a productive infinitive confuses a citation with evidence for a form. [2]
% mdparagraph:P043:end

\end{minipage}\par\addvspace{4pt}

\reserve{7\baselineskip}
% mdparagraph:P044:begin
\section*{M. Knowledge, ability, necessity and willingness}
% mdparagraph:P044:end

\par\begin{minipage}{\linewidth}\raggedright
% item: NEW-U06-M01
% key: KEY-NEW-U06-M01
% mdparagraph:P045:begin
\textbf{M01.}
% mdparagraph:P045:end

% mdparagraph:P046:begin
\textbf{Ic cann swimman. Wē cunnon singan.}
% mdparagraph:P046:end

% mdparagraph:P047:begin
The target is learned ability: “I know how to swim; we know how to sing.” \textbf{Cann} is present indicative first singular; \textbf{cunnon} is first plural here. \textbf{Swimman} and \textbf{singan} are plain infinitives. English “can” is also defensible if the explanation preserves the skill sense. A second finite verb or an inflected infinitive after the modal misses the construction. [5]
% mdparagraph:P047:end

\end{minipage}\par\addvspace{4pt}

\par\begin{minipage}{\linewidth}\raggedright
% item: NEW-U06-M02
% key: KEY-NEW-U06-M02
% mdparagraph:P048:begin
\textbf{M02.}
% mdparagraph:P048:end

% mdparagraph:P049:begin
\textbf{Hē mæg ðæt wæter beran. Hīe magon ðǣm męnn helpan.}
% mdparagraph:P049:end

% mdparagraph:P050:begin
The target is present capability: he can carry the water; they can help the man. \textbf{Mæg/\allowbreak{}magon} agree with the singular/\allowbreak{}plural subjects. \textbf{Ðæt wæter} is the accusative thing carried; \textbf{ðǣm męnn} is the dative person helped. A different defensible order of those constituents is acceptable, but changing the helped person to an accusative loses the taught government. The infinitives remain unchanged by subject number. [5]
% mdparagraph:P050:end

\end{minipage}\par\addvspace{4pt}

\par\begin{minipage}{\linewidth}\raggedright
% item: NEW-U06-M03
% key: KEY-NEW-U06-M03
% mdparagraph:P051:begin
\textbf{M03.}
% mdparagraph:P051:end

% mdparagraph:P052:begin
\textbf{Ic ðearf gān. Wē ðorfton gān.}
% mdparagraph:P052:end

% mdparagraph:P053:begin
The first clause expresses my present need; the second our past need. \textbf{Ðearf} is first-singular present and \textbf{ðorfton} first-plural preterite in these clauses. \textbf{Gān} remains plain in both. “We had a need to go” preserves the second meaning; a merely hypothetical “we would need to go” adds an unsupported condition. Keep the past \textbf{o} distinct from present-plural \textbf{u}. [5]
% mdparagraph:P053:end

\end{minipage}\par\addvspace{4pt}

\par\begin{minipage}{\linewidth}\raggedright
% item: NEW-U06-M04
% key: KEY-NEW-U06-M04
% mdparagraph:P054:begin
\textbf{M04.}
% mdparagraph:P054:end

% mdparagraph:P055:begin
\textbf{Ðū dearst gān. Ðū dorstest gān.}
% mdparagraph:P055:end

% mdparagraph:P056:begin
Both subjects are one addressee. \textbf{Dearst} is second-singular present indicative, \textbf{dorstest} second-singular preterite. The target contrast is “dare” against actual past “dared.” A form belonging to first/\allowbreak{}third singular misses the explicit second person; the plain infinitive \textbf{gān} does not change with tense. [5]
% mdparagraph:P056:end

\end{minipage}\par\addvspace{4pt}

\par\begin{minipage}{\linewidth}\raggedright
% item: NEW-U06-M05
% key: KEY-NEW-U06-M05
% mdparagraph:P057:begin
\textbf{M05.}
% mdparagraph:P057:end

% mdparagraph:P058:begin
\textbf{Wē sculon ðæt weorc dōn. Hīe sceoldon ðæt weorc dōn.}
% mdparagraph:P058:end

% mdparagraph:P059:begin
The first clause states our obligation now; the second their obligation in the past. \textbf{Sculon} is first plural here, \textbf{sceoldon} third plural; \textbf{dōn} is the plain infinitive with accusative \textbf{ðæt weorc}. “Were obliged to do” is a good rendering of the past. Other taught wording for performing the work can preserve the meaning, but the requested obligation verb remains \textbf{sculan}. Predictive “will do” or hypothetical “would do” would change this target. The documented \textbf{scoldon} variant remains linguistically valid, though \textbf{sceoldon} is the selected form. [5]
% mdparagraph:P059:end

\end{minipage}\par\addvspace{4pt}

\par\begin{minipage}{\linewidth}\raggedright
% item: NEW-U06-M06
% key: KEY-NEW-U06-M06
% mdparagraph:P060:begin
\textbf{M06.}
% mdparagraph:P060:end

% mdparagraph:P061:begin
\textbf{Ðū mōst ðone męte etan. Ðū mōstest ðone męte etan.}
% mdparagraph:P061:end

% mdparagraph:P062:begin
The first sentence gives one person present permission to eat the food; the second gives that person past permission. \textbf{Mōst} and \textbf{mōstest} are second singular, with long \textbf{ō}. \textbf{Etan} takes the accusative food object in both. “You may eat” is suitable for the first sentence. English “must eat” selects a different sense and does not satisfy the stated permission context. [5]
% mdparagraph:P062:end

\end{minipage}\par\addvspace{4pt}

\par\begin{minipage}{\linewidth}\raggedright
% item: NEW-U06-M07
% key: KEY-NEW-U06-M07
% mdparagraph:P063:begin
\textbf{M07.}
% mdparagraph:P063:end

% mdparagraph:P064:begin
\textbf{Hē wille ðǣm fuglum wæter giefan. Hīe woldon ðǣm fuglum wæter giefan.}
% mdparagraph:P064:end

% mdparagraph:P065:begin
The target is his present desire and their past desire to give water to the birds. \textbf{Wille} is third-singular present here, \textbf{woldon} third-plural preterite. \textbf{Wæter} is accusative; \textbf{ðǣm fuglum} is the dative recipient. The separately taught \textbf{sęllan} can also express giving in this frame. Keep the volitional verb and its tense; neither sentence by itself asserts that the giving happened. Adding an ordinary third-singular ending to \textbf{wille} is not the selected paradigm. [5]
% mdparagraph:P065:end

\end{minipage}\par\addvspace{4pt}

\par\begin{minipage}{\linewidth}\raggedright
% item: NEW-U06-M08
% key: KEY-NEW-U06-M08
% mdparagraph:P066:begin
\textbf{M08.}
% mdparagraph:P066:end

% mdparagraph:P067:begin
\textbf{Ic ne dorste gān. Wē ne magon ðæt scip beran.}
% mdparagraph:P067:end

% mdparagraph:P068:begin
The first sentence denies past daring; the second denies present ability. Put \textbf{ne} before the finite modal. \textbf{Dorste} is first-singular preterite here, \textbf{magon} first-plural present. \textbf{Ðæt scip} is the accusative thing carried. A different defensible order may preserve the relations, but negating only the embedded action would not express these same modal meanings. No new contracted negative form is needed. [5]
% mdparagraph:P068:end

\end{minipage}\par\addvspace{4pt}

\par\begin{minipage}{\linewidth}\raggedright
% item: NEW-U06-M09
% key: KEY-NEW-U06-M09
% mdparagraph:P069:begin
\textbf{M09.}
% mdparagraph:P069:end

% mdparagraph:P070:begin
\textbf{Wit cunnon swimman. Git sculon unc helpan.}
% mdparagraph:P070:end

% mdparagraph:P071:begin
\textbf{Wit} means “we two,” \textbf{git} “you two,” and \textbf{unc} “us two.” The first sentence expresses learned ability; the second obliges the two addressees to help the two speakers. Dual subjects take ordinary plural verbs, hence \textbf{cunnon} and \textbf{sculon}. \textbf{Unc} is dative because \textbf{helpan} governs the helped person in that case. Ordinary plural pronouns would lose the expressly requested two-person reference; a separate dual verb ending is not supplied.
% mdparagraph:P071:end

\end{minipage}\par\addvspace{4pt}

\par\begin{minipage}{\linewidth}\raggedright
% item: NEW-U06-M10
% key: KEY-NEW-U06-M10
% mdparagraph:P072:begin
\textbf{M10.}
% mdparagraph:P072:end

% mdparagraph:P073:begin
\textbf{Hē wolde gān. Hē ēode.}
% mdparagraph:P073:end

% mdparagraph:P074:begin
The first clause says that he wanted or intended to go; the second says that he went. \textbf{Wolde} is preterite indicative third singular with plain \textbf{gān}; \textbf{ēode} is preterite indicative third singular of \textbf{gān}. A wish can remain unfulfilled, so the first statement alone does not entail the second. Defensible wording must preserve that difference; treating \textbf{wolde} as automatically reporting an accomplished action misses the point.
% mdparagraph:P074:end

\end{minipage}\par\addvspace{4pt}

\par\begin{minipage}{\linewidth}\raggedright
% item: NEW-U06-M11
% key: KEY-NEW-U06-M11
% mdparagraph:P075:begin
\textbf{M11.}
% mdparagraph:P075:end

% mdparagraph:P076:begin
\textbf{Ic cann swimman:} “I know how to swim.” \textbf{Ic mæg swimman:} “I can swim,” in the stated sense of present capability.
% mdparagraph:P076:end

% mdparagraph:P077:begin
\textbf{Cann} and \textbf{mæg} are the finite present indicative verbs, first singular here. \textbf{Swimman} is the plain infinitive in both and names the action. The first sentence emphasizes learned skill; the second emphasizes ability or capacity in the relevant circumstances. English “can” can render either, provided the distinction is explained. This exercise does not establish a universal one-to-one mapping between English auxiliaries and Old English verbs. [5]
% mdparagraph:P077:end

\end{minipage}\par\addvspace{4pt}

\par\begin{minipage}{\linewidth}\raggedright
% item: NEW-U06-M12
% key: KEY-NEW-U06-M12
% mdparagraph:P078:begin
\textbf{M12.}
% mdparagraph:P078:end

% mdparagraph:P079:begin
\textbf{Ðæt scip bið gearu. Ðæt scip is gearu.}
% mdparagraph:P079:end

% mdparagraph:P080:begin
In the given expected-later situation, the first means “The ship will be ready.” The second gives the requested ordinary present comparison, “The ship is ready.” \textbf{Bið} and \textbf{is} are third-singular indicative forms. \textbf{Gearu} is the strong nominative neuter singular predicate adjective agreeing with \textbf{scip}; the subject’s article does not make a predicate adjective weak.
% mdparagraph:P080:end

% mdparagraph:P081:begin
The bēo series can also express habitual or general truths, so “future” is the contextual interpretation here, not the meaning of every \textbf{bið}. An answer using the eom series for both sentences may be defensible in some future context but would miss the first sentence’s expressly requested series. [2]
% mdparagraph:P081:end

\end{minipage}\par\addvspace{4pt}

\par\begin{minipage}{\linewidth}\raggedright
% item: NEW-U06-M13
% key: KEY-NEW-U06-M13
% mdparagraph:P082:begin
\textbf{M13.}
% mdparagraph:P082:end

% mdparagraph:P083:begin
\textbf{Ic wāt ðā ðing. Hē wiste ðā ðing.}
% mdparagraph:P083:end

% mdparagraph:P084:begin
The target is present “I know” and past “he knew,” each with the accusative nominal object “those things.” \textbf{Ðā} is the pointing demonstrative; \textbf{ðing} has its explicitly taught fact/\allowbreak{}matter sense here. The noun’s plural form happens to match its nominative, but its relation to \textbf{witan} is accusative.
% mdparagraph:P084:end

% mdparagraph:P085:begin
This construction names what is known. \textbf{Cunnan} with a plain infinitive, as in M01, expresses knowing how to perform an action. The distinction is between the assigned constructions and senses, not a claim that \textbf{cunnan} has no other uses anywhere in Old English. Using it instead of the requested \textbf{witan} would miss this task’s target. [6]
% mdparagraph:P085:end

\end{minipage}\par\addvspace{4pt}

\reserve{7\baselineskip}
% mdparagraph:P086:begin
\section*{I. Commands}
% mdparagraph:P086:end

\par\begin{minipage}{\linewidth}\raggedright
% item: NEW-U06-I01
% key: KEY-NEW-U06-I01
% mdparagraph:P087:begin
\textbf{I01.}
% mdparagraph:P087:end

% mdparagraph:P088:begin
\textbf{Target: the complete second-person command pairs.}
% mdparagraph:P088:end

% mdtable:T13:begin
% course-table-style: 2026-10
% course-panel:T13-P01:begin
\par\begingroup
\fontsize{10.5}{13.2}\selectfont
\setlength{\LTleft}{0pt}\setlength{\LTright}{\fill}
\noindent
\begin{tabular}{@{}>{\raggedright\arraybackslash}p{30.902mm}>{\raggedright\arraybackslash}p{28.221mm}>{\raggedright\arraybackslash}p{28.221mm}>{\raggedright\arraybackslash}p{28.221mm}@{}}
\toprule
\textbf{Person and number} & \textbf{hīeran} & \textbf{nęrian} & \textbf{gefręmman} \\
\midrule
2 sg. & hīer & nęre & gefręme \\
2 pl. & hīerað & nęriað & gefręmmað \\
\bottomrule
\end{tabular}
\par\endgroup
% course-panel:T13-P01:end
% mdtable:T13:end
% mdparagraph:P089:begin
\textbf{Hīer} has no final \textbf{-e}; \textbf{nęre} and \textbf{gefręme} retain it. The plural restores the selected \textbf{-iað} or geminated stem. Merely removing the final \textbf{-n} of each infinitive will not give these singulars. [3]
% mdparagraph:P089:end

\end{minipage}\par\addvspace{4pt}

\par\begin{minipage}{\linewidth}\raggedright
% item: NEW-U06-I02
% key: KEY-NEW-U06-I02
% mdparagraph:P090:begin
\textbf{I02.}
% mdparagraph:P090:end

% mdparagraph:P091:begin
\textbf{Target: the complete second-person command pairs.}
% mdparagraph:P091:end

% mdtable:T14:begin
% course-table-style: 2026-10
% course-panel:T14-P01:begin
\par\begingroup
\fontsize{10.5}{13.2}\selectfont
\setlength{\LTleft}{0pt}\setlength{\LTright}{\fill}
\noindent
\begin{tabular}{@{}>{\raggedright\arraybackslash}p{24.232mm}>{\raggedright\arraybackslash}p{22.13mm}>{\raggedright\arraybackslash}p{22.13mm}>{\raggedright\arraybackslash}p{22.13mm}>{\raggedright\arraybackslash}p{22.13mm}@{}}
\toprule
\textbf{Person and number} & \textbf{cyssan} & \textbf{gefyllan} & \textbf{sęttan} & \textbf{gierwan} \\
\midrule
2 sg. & cys & gefyl & sęte & giere \\
2 pl. & cyssað & gefyllað & sęttað & gierwað \\
\bottomrule
\end{tabular}
\par\endgroup
% course-panel:T14-P01:end
\par\addvspace{5pt}
% course-panel:T14-P02:begin
\par\begingroup
\fontsize{10.5}{13.2}\selectfont
\setlength{\LTleft}{0pt}\setlength{\LTright}{\fill}
\noindent
\begin{tabular}{@{}>{\raggedright\arraybackslash}p{63.342mm}>{\raggedright\arraybackslash}p{57.846mm}@{}}
\toprule
\textbf{Person and number} & \textbf{sęllan} \\
\midrule
2 sg. & sęle \\
2 pl. & sęllað \\
\bottomrule
\end{tabular}
\par\endgroup
% course-panel:T14-P02:end
% mdtable:T14:end
% mdparagraph:P092:begin
The selected heavy-stem singulars \textbf{cys} and \textbf{gefyl} simplify final geminates. The short-stem forms \textbf{sęte} and \textbf{sęle} retain \textbf{-e} while losing gemination; their plurals retain \textbf{tt/\allowbreak{}ll}. \textbf{Giere} lacks the \textbf{w} retained in \textbf{gierwað}. The singular pattern is not the same for every doubled consonant. Documented unsimplified spellings need not be impossible Old English, but these are the requested selections. [3]
% mdparagraph:P092:end

\end{minipage}\par\addvspace{4pt}

\par\begin{minipage}{\linewidth}\raggedright
% item: NEW-U06-I03
% key: KEY-NEW-U06-I03
% mdparagraph:P093:begin
\textbf{I03.}
% mdparagraph:P093:end

% mdparagraph:P094:begin
\textbf{Target: the complete second-person command pairs.}
% mdparagraph:P094:end

% mdtable:T15:begin
% course-table-style: 2026-10
% course-panel:T15-P01:begin
\par\begingroup
\fontsize{10.5}{13.2}\selectfont
\setlength{\LTleft}{0pt}\setlength{\LTright}{\fill}
\noindent
\begin{tabular}{@{}>{\raggedright\arraybackslash}p{30.902mm}>{\raggedright\arraybackslash}p{28.221mm}>{\raggedright\arraybackslash}p{28.221mm}>{\raggedright\arraybackslash}p{28.221mm}@{}}
\toprule
\textbf{Person and number} & \textbf{lufian} & \textbf{habban} & \textbf{libban} \\
\midrule
2 sg. & lufa & hafa & leofa \\
2 pl. & lufiað & habbað & libbað \\
\bottomrule
\end{tabular}
\par\endgroup
% course-panel:T15-P01:end
% mdtable:T15:end
% mdparagraph:P095:begin
Weak II has \textbf{lufa/\allowbreak{}lufiað}. The weak III survivors require their own pairs: \textbf{hafa/\allowbreak{}habbað} and \textbf{leofa/\allowbreak{}libbað}. Do not impose the weak-II singular ending on an unchanged weak-III infinitive stem. Cook also gives the documented \textbf{lifiað} plural alternative; the selected plural here is \textbf{libbað}. [3]
% mdparagraph:P095:end

\end{minipage}\par\addvspace{4pt}

\par\begin{minipage}{\linewidth}\raggedright
% item: NEW-U06-I04
% key: KEY-NEW-U06-I04
% mdparagraph:P096:begin
\textbf{I04.}
% mdparagraph:P096:end

% mdparagraph:P097:begin
\textbf{Target: the complete second-person command pairs.}
% mdparagraph:P097:end

% mdtable:T16:begin
% course-table-style: 2026-10
% course-panel:T16-P01:begin
\par\begingroup
\fontsize{10.5}{13.2}\selectfont
\setlength{\LTleft}{0pt}\setlength{\LTright}{\fill}
\noindent
\begin{tabular}{@{}>{\raggedright\arraybackslash}p{39.459mm}>{\raggedright\arraybackslash}p{39.459mm}>{\raggedright\arraybackslash}p{39.459mm}@{}}
\toprule
\textbf{Person and number} & \textbf{sęcgan} & \textbf{hycgan} \\
\midrule
2 sg. & saga & hoga \\
2 pl. & sęcgað & hycgað \\
\bottomrule
\end{tabular}
\par\endgroup
% course-panel:T16-P01:end
% mdtable:T16:end
% mdparagraph:P098:begin
The singulars \textbf{saga} and \textbf{hoga} differ sharply from the selected plural stems. Cook also records \textbf{sęge} beside \textbf{saga}, and optional \textbf{e} in the plural spellings. These are documented alternatives; the table supplies the requested course forms. Knowing a command form does not by itself establish a new object or complement construction. [3]
% mdparagraph:P098:end

\end{minipage}\par\addvspace{4pt}

\par\begin{minipage}{\linewidth}\raggedright
% item: NEW-U06-I05
% key: KEY-NEW-U06-I05
% mdparagraph:P099:begin
\textbf{I05.}
% mdparagraph:P099:end

% mdparagraph:P100:begin
\textbf{Target: the complete second-person command pairs.}
% mdparagraph:P100:end

% mdtable:T17:begin
% course-table-style: 2026-10
% course-panel:T17-P01:begin
\par\begingroup
\fontsize{10.5}{13.2}\selectfont
\setlength{\LTleft}{0pt}\setlength{\LTright}{\fill}
\noindent
\begin{tabular}{@{}>{\raggedright\arraybackslash}p{24.232mm}>{\raggedright\arraybackslash}p{22.13mm}>{\raggedright\arraybackslash}p{22.13mm}>{\raggedright\arraybackslash}p{22.13mm}>{\raggedright\arraybackslash}p{22.13mm}@{}}
\toprule
\textbf{Person and number} & \textbf{drīfan} & \textbf{bēodan} & \textbf{bindan} & \textbf{helpan} \\
\midrule
2 sg. & drīf & bēod & bind & help \\
2 pl. & drīfað & bēodað & bindað & helpað \\
\bottomrule
\end{tabular}
\par\endgroup
% course-panel:T17-P01:end
\par\addvspace{5pt}
% course-panel:T17-P02:begin
\par\begingroup
\fontsize{10.5}{13.2}\selectfont
\setlength{\LTleft}{0pt}\setlength{\LTright}{\fill}
\noindent
\begin{tabular}{@{}>{\raggedright\arraybackslash}p{63.342mm}>{\raggedright\arraybackslash}p{57.846mm}@{}}
\toprule
\textbf{Person and number} & \textbf{drincan} \\
\midrule
2 sg. & drinc \\
2 pl. & drincað \\
\bottomrule
\end{tabular}
\par\endgroup
% course-panel:T17-P02:end
% mdtable:T17:end
% mdparagraph:P101:begin
Use the present stems shown here. Keep long \textbf{ī} in \textbf{drīf} and long \textbf{ēo} in \textbf{bēod}; \textbf{bind, help} and \textbf{drinc} have their selected short vowels. A changed strong-past stem is not a command stem. The plural \textbf{-að} addresses more than one person. [3]
% mdparagraph:P101:end

\end{minipage}\par\addvspace{4pt}

\par\begin{minipage}{\linewidth}\raggedright
% item: NEW-U06-I06
% key: KEY-NEW-U06-I06
% mdparagraph:P102:begin
\textbf{I06.}
% mdparagraph:P102:end

% mdparagraph:P103:begin
\textbf{Target: the complete second-person command pairs.}
% mdparagraph:P103:end

% mdtable:T18:begin
% course-table-style: 2026-10
% course-panel:T18-P01:begin
\par\begingroup
\fontsize{10.5}{13.2}\selectfont
\setlength{\LTleft}{0pt}\setlength{\LTright}{\fill}
\noindent
\begin{tabular}{@{}>{\raggedright\arraybackslash}p{39.459mm}>{\raggedright\arraybackslash}p{39.459mm}>{\raggedright\arraybackslash}p{39.459mm}@{}}
\toprule
\textbf{Person and number} & \textbf{wrēon} & \textbf{tēon} \\
\midrule
2 sg. & wrīh & tēoh \\
2 pl. & wrēoð & tēoð \\
\bottomrule
\end{tabular}
\par\endgroup
% course-panel:T18-P01:end
% mdtable:T18:end
% mdparagraph:P104:begin
The singular commands restore \textbf{h}, but their vowels differ: \textbf{wrīh} has long \textbf{ī}, \textbf{tēoh} long \textbf{ēo}. The plurals remain contracted. Do not transfer either singular vowel to the other verb or append the singular \textbf{h} to the selected plural. [3]
% mdparagraph:P104:end

\end{minipage}\par\addvspace{4pt}

\par\begin{minipage}{\linewidth}\raggedright
% item: NEW-U06-I07
% key: KEY-NEW-U06-I07
% mdparagraph:P105:begin
\textbf{I07.}
% mdparagraph:P105:end

% mdparagraph:P106:begin
\textbf{Target: the complete second-person command pairs.}
% mdparagraph:P106:end

% mdtable:T19:begin
% course-table-style: 2026-10
% course-panel:T19-P01:begin
\par\begingroup
\fontsize{10.5}{13.2}\selectfont
\setlength{\LTleft}{0pt}\setlength{\LTright}{\fill}
\noindent
\begin{tabular}{@{}>{\raggedright\arraybackslash}p{30.902mm}>{\raggedright\arraybackslash}p{28.221mm}>{\raggedright\arraybackslash}p{28.221mm}>{\raggedright\arraybackslash}p{28.221mm}@{}}
\toprule
\textbf{Person and number} & \textbf{beran} & \textbf{niman} & \textbf{cuman} \\
\midrule
2 sg. & ber & nim & cum \\
2 pl. & berað & nimað & cumað \\
\bottomrule
\end{tabular}
\par\endgroup
% course-panel:T19-P01:end
% mdtable:T19:end
% mdparagraph:P107:begin
Use \textbf{ber, nim, cum}, without importing the vowel changes of present indicative second/\allowbreak{}third singular. All three plural commands use \textbf{-að} on the selected present stem. Replacing a command by a personal indicative form changes the requested mood. [3]
% mdparagraph:P107:end

\end{minipage}\par\addvspace{4pt}

\par\begin{minipage}{\linewidth}\raggedright
% item: NEW-U06-I08
% key: KEY-NEW-U06-I08
% mdparagraph:P108:begin
\textbf{I08.}
% mdparagraph:P108:end

% mdparagraph:P109:begin
\textbf{Target: the complete second-person command pairs.}
% mdparagraph:P109:end

% mdtable:T20:begin
% course-table-style: 2026-10
% course-panel:T20-P01:begin
\par\begingroup
\fontsize{10.5}{13.2}\selectfont
\setlength{\LTleft}{0pt}\setlength{\LTright}{\fill}
\noindent
\begin{tabular}{@{}>{\raggedright\arraybackslash}p{24.232mm}>{\raggedright\arraybackslash}p{22.13mm}>{\raggedright\arraybackslash}p{22.13mm}>{\raggedright\arraybackslash}p{22.13mm}>{\raggedright\arraybackslash}p{22.13mm}@{}}
\toprule
\textbf{Person and number} & \textbf{giefan} & \textbf{biddan} & \textbf{sittan} & \textbf{licgan} \\
\midrule
2 sg. & gief & bide & site & lige \\
2 pl. & giefað & biddað & sittað & licgað \\
\bottomrule
\end{tabular}
\par\endgroup
% course-panel:T20-P01:end
\par\addvspace{5pt}
% course-panel:T20-P02:begin
\par\begingroup
\fontsize{10.5}{13.2}\selectfont
\setlength{\LTleft}{0pt}\setlength{\LTright}{\fill}
\noindent
\begin{tabular}{@{}>{\raggedright\arraybackslash}p{63.342mm}>{\raggedright\arraybackslash}p{57.846mm}@{}}
\toprule
\textbf{Person and number} & \textbf{etan} \\
\midrule
2 sg. & et \\
2 pl. & etað \\
\bottomrule
\end{tabular}
\par\endgroup
% course-panel:T20-P02:end
% mdtable:T20:end
% mdparagraph:P110:begin
The j-presents have \textbf{bide, site, lige} in the singular, with single consonants and final \textbf{-e}; their plurals retain the selected stems. \textbf{Gief} is the ordinary stem command. \textbf{Et} has neither the changed dental present-singular stem nor the long vowel of the past. [3]
% mdparagraph:P110:end

\end{minipage}\par\addvspace{4pt}

\par\begin{minipage}{\linewidth}\raggedright
% item: NEW-U06-I09
% key: KEY-NEW-U06-I09
% mdparagraph:P111:begin
\textbf{I09.}
% mdparagraph:P111:end

% mdparagraph:P112:begin
\textbf{Target: the complete second-person command pairs.}
% mdparagraph:P112:end

% mdtable:T21:begin
% course-table-style: 2026-10
% course-panel:T21-P01:begin
\par\begingroup
\fontsize{10.5}{13.2}\selectfont
\setlength{\LTleft}{0pt}\setlength{\LTright}{\fill}
\noindent
\begin{tabular}{@{}>{\raggedright\arraybackslash}p{30.902mm}>{\raggedright\arraybackslash}p{28.221mm}>{\raggedright\arraybackslash}p{28.221mm}>{\raggedright\arraybackslash}p{28.221mm}@{}}
\toprule
\textbf{Person and number} & \textbf{sēon} & \textbf{slēan} & \textbf{fōn} \\
\midrule
2 sg. & seoh & sleah & fōh \\
2 pl. & sēoð & slēað & fōð \\
\bottomrule
\end{tabular}
\par\endgroup
% course-panel:T21-P01:end
% mdtable:T21:end
% mdparagraph:P113:begin
Singular \textbf{seoh} and \textbf{sleah} have short diphthongs; \textbf{fōh} has long \textbf{ō}. The restored \textbf{h} belongs to the singular command. The plurals are \textbf{sēoð, slēað, fōð}, preserving their own quantities and contractions. [3]
% mdparagraph:P113:end

\end{minipage}\par\addvspace{4pt}

\par\begin{minipage}{\linewidth}\raggedright
% item: NEW-U06-I10
% key: KEY-NEW-U06-I10
% mdparagraph:P114:begin
\textbf{I10.}
% mdparagraph:P114:end

% mdparagraph:P115:begin
\textbf{Target: the complete second-person command pairs.}
% mdparagraph:P115:end

% mdtable:T22:begin
% course-table-style: 2026-10
% course-panel:T22-P01:begin
\par\begingroup
\fontsize{10.5}{13.2}\selectfont
\setlength{\LTleft}{0pt}\setlength{\LTright}{\fill}
\noindent
\begin{tabular}{@{}>{\raggedright\arraybackslash}p{30.902mm}>{\raggedright\arraybackslash}p{28.221mm}>{\raggedright\arraybackslash}p{28.221mm}>{\raggedright\arraybackslash}p{28.221mm}@{}}
\toprule
\textbf{Person and number} & \textbf{hębban} & \textbf{scieppan} & \textbf{swęrian} \\
\midrule
2 sg. & hęfe & sciepe & swęre \\
2 pl. & hębbað & scieppað & swęriað \\
\bottomrule
\end{tabular}
\par\endgroup
% course-panel:T22-P01:end
% mdtable:T22:end
% mdparagraph:P116:begin
These singular commands use \textbf{-e}. Note \textbf{f/\allowbreak{}bb} in \textbf{hęfe/\allowbreak{}hębbað}, single \textbf{p/\allowbreak{}pp} in \textbf{sciepe/\allowbreak{}scieppað}, and \textbf{swęre/\allowbreak{}swęriað}. A single rule of cutting off an infinitive ending misses these stem alternations. [3]
% mdparagraph:P116:end

\end{minipage}\par\addvspace{4pt}

\par\begin{minipage}{\linewidth}\raggedright
% item: NEW-U06-I11
% key: KEY-NEW-U06-I11
% mdparagraph:P117:begin
\textbf{I11.}
% mdparagraph:P117:end

% mdparagraph:P118:begin
\textbf{Target: the complete second-person command pairs.}
% mdparagraph:P118:end

% mdtable:T23:begin
% course-table-style: 2026-10
% course-panel:T23-P01:begin
\par\begingroup
\fontsize{10.5}{13.2}\selectfont
\setlength{\LTleft}{0pt}\setlength{\LTright}{\fill}
\noindent
\begin{tabular}{@{}>{\raggedright\arraybackslash}p{30.902mm}>{\raggedright\arraybackslash}p{28.221mm}>{\raggedright\arraybackslash}p{28.221mm}>{\raggedright\arraybackslash}p{28.221mm}@{}}
\toprule
\textbf{Person and number} & \textbf{faran} & \textbf{standan} & \textbf{healdan} \\
\midrule
2 sg. & far & stand & heald \\
2 pl. & farað & standað & healdað \\
\bottomrule
\end{tabular}
\par\endgroup
% course-panel:T23-P01:end
% mdtable:T23:end
% mdparagraph:P119:begin
The command stems are \textbf{far, stand, heald}. Do not substitute a past stem or the altered stem of an indicative singular. The paired plural changes the number of addressees, not the tense. [3]
% mdparagraph:P119:end

\end{minipage}\par\addvspace{4pt}

\par\begin{minipage}{\linewidth}\raggedright
% item: NEW-U06-I12
% key: KEY-NEW-U06-I12
% mdparagraph:P120:begin
\textbf{I12.}
% mdparagraph:P120:end

% mdparagraph:P121:begin
\textbf{Target: the complete second-person command pairs.}
% mdparagraph:P121:end

% mdtable:T24:begin
% course-table-style: 2026-10
% course-panel:T24-P01:begin
\par\begingroup
\fontsize{10.5}{13.2}\selectfont
\setlength{\LTleft}{0pt}\setlength{\LTright}{\fill}
\noindent
\begin{tabular}{@{}>{\raggedright\arraybackslash}p{30.902mm}>{\raggedright\arraybackslash}p{28.221mm}>{\raggedright\arraybackslash}p{28.221mm}>{\raggedright\arraybackslash}p{28.221mm}@{}}
\toprule
\textbf{Person and number} & \textbf{weaxan} & \textbf{hātan} & \textbf{lǣtan} \\
\midrule
2 sg. & weax & hāt & lǣt \\
2 pl. & weaxað & hātað & lǣtað \\
\bottomrule
\end{tabular}
\par\endgroup
% course-panel:T24-P01:end
% mdtable:T24:end
% mdparagraph:P122:begin
Preserve \textbf{ea} in \textbf{weax}, long \textbf{ā} in \textbf{hāt}, and long \textbf{ǣ} in \textbf{lǣt}. The commands use these selected present stems rather than the changed or contracted forms of indicative second/\allowbreak{}third singular. [3]
% mdparagraph:P122:end

\end{minipage}\par\addvspace{4pt}

\par\begin{minipage}{\linewidth}\raggedright
% item: NEW-U06-I13
% key: KEY-NEW-U06-I13
% mdparagraph:P123:begin
\textbf{I13.}
% mdparagraph:P123:end

% mdparagraph:P124:begin
\textbf{Target: the complete second-person command pairs.}
% mdparagraph:P124:end

% mdtable:T25:begin
% course-table-style: 2026-10
% course-panel:T25-P01:begin
\par\begingroup
\fontsize{10.5}{13.2}\selectfont
\setlength{\LTleft}{0pt}\setlength{\LTright}{\fill}
\noindent
\begin{tabular}{@{}>{\raggedright\arraybackslash}p{24.232mm}>{\raggedright\arraybackslash}p{22.13mm}>{\raggedright\arraybackslash}p{22.13mm}>{\raggedright\arraybackslash}p{22.13mm}>{\raggedright\arraybackslash}p{22.13mm}@{}}
\toprule
\textbf{Person and number} & \textbf{dōn} & \textbf{gān} & \textbf{wesan} & \textbf{bēon} \\
\midrule
2 sg. & dō & gā & wes & bēo \\
2 pl. & dōð & gāð & wesað & bēoð \\
\bottomrule
\end{tabular}
\par\endgroup
% course-panel:T25-P01:end
% mdtable:T25:end
% mdparagraph:P125:begin
Both be-command series are available. \textbf{Dō, gā, wes, bēo} address one person; their paired plurals address two or more. \textbf{Bēo} here is a command, not the first-singular indicative merely because the two cells share a spelling. No positive imperative of \textbf{willan} or regular modal command pair is required. [1,3]
% mdparagraph:P125:end

\end{minipage}\par\addvspace{4pt}

\par\begin{minipage}{\linewidth}\raggedright
% item: NEW-U06-I14
% key: KEY-NEW-U06-I14
% mdparagraph:P126:begin
\textbf{I14.}
% mdparagraph:P126:end

% mdparagraph:P127:begin
\textbf{One addressee:} \textbf{Nim ðone męte and ber hine.}
% mdparagraph:P127:end

% mdparagraph:P128:begin
\textbf{Several addressees:} \textbf{Nimað ðone męte and berað hine.}
% mdparagraph:P128:end

% mdparagraph:P129:begin
Both commands mean “Take the food and carry it.” \textbf{Hine} is masculine accusative singular, referring to \textbf{męte}, also masculine. The object keeps its number when the addressees change. A defensible word order can vary, but a neuter pronoun chosen from English “it” would miss the Old English noun’s gender. Both coordinated verbs must have the requested command number.
% mdparagraph:P129:end

\end{minipage}\par\addvspace{4pt}

\par\begin{minipage}{\linewidth}\raggedright
% item: NEW-U06-I15
% key: KEY-NEW-U06-I15
% mdparagraph:P130:begin
\textbf{I15.}
% mdparagraph:P130:end

% mdparagraph:P131:begin
\textbf{One addressee:} \textbf{Help mē. Gief mē ðæt wæter.}
% mdparagraph:P131:end

% mdparagraph:P132:begin
\textbf{Two addressees:} \textbf{Helpað mē. Giefað mē ðæt wæter.}
% mdparagraph:P132:end

% mdparagraph:P133:begin
In both versions, \textbf{mē} is dative: the person helped and the recipient of the water. \textbf{Ðæt wæter} is the accusative gift. Exactly two addressees take the ordinary plural imperative; the commands have no separate dual ending. The taught giving verb \textbf{sęllan} can preserve the gift relation with its own command forms. Changing \textbf{mē} to an object form selected without regard to the verb’s government misses the target.
% mdparagraph:P133:end

\end{minipage}\par\addvspace{4pt}

\par\begin{minipage}{\linewidth}\raggedright
% item: NEW-U06-I16
% key: KEY-NEW-U06-I16
% mdparagraph:P134:begin
\textbf{I16.}
% mdparagraph:P134:end

% mdparagraph:P135:begin
\textbf{Ne drincað ðæt wæter. Sittað and etað.}
% mdparagraph:P135:end

% mdparagraph:P136:begin
All three verbs are second-person plural imperatives. \textbf{Ne} negates the first command: “Do not drink the water.” The next two positively command sitting and eating. \textbf{Ðæt wæter} is the accusative object of drinking; eating may stand without an expressed food object. Singular commands would address the wrong number of people. The negative construction does not require an additional auxiliary. [7]
% mdparagraph:P136:end

\end{minipage}\par\addvspace{4pt}

\par\begin{minipage}{\linewidth}\raggedright
% item: NEW-U06-I17
% key: KEY-NEW-U06-I17
% mdparagraph:P137:begin
\textbf{I17.}
% mdparagraph:P137:end

% mdparagraph:P138:begin
\textbf{Se mann cwæð: ‘Gā and dō ðæt weorc wel.’}
% mdparagraph:P138:end

% mdparagraph:P139:begin
\textbf{Cwæð} reports the man’s past speech. Within the quotation, \textbf{gā} and \textbf{dō} are singular commands to his one companion; \textbf{ðæt weorc} is the accusative work to be done and \textbf{wel} means “well.” The command remains a direct utterance, so it is not converted into a past indicative merely because the reporting verb is past. A defensible placement of \textbf{wel} is acceptable if its relation to doing remains clear.
% mdparagraph:P139:end

\end{minipage}\par\addvspace{4pt}

\reserve{7\baselineskip}
% mdparagraph:P140:begin
\section*{N. Infinitives}
% mdparagraph:P140:end

\par\begin{minipage}{\linewidth}\raggedright
% item: NEW-U06-N01
% key: KEY-NEW-U06-N01
% mdparagraph:P141:begin
\textbf{N01.}
% mdparagraph:P141:end

% mdparagraph:P142:begin
\textbf{Target: the plain and selected inflected infinitives.}
% mdparagraph:P142:end

% mdtable:T26:begin
\par\begingroup
\fontsize{10.5}{13.2}\selectfont
\setlength{\LTleft}{0pt}\setlength{\LTright}{\fill}
\noindent
\begin{tabular}{@{}>{\raggedright\arraybackslash}p{67.8mm}>{\raggedright\arraybackslash}p{53.39mm}@{}}
\toprule
\textbf{Plain infinitive} & \textbf{After tō} \\
\midrule
bindan & tō bindanne \\
beran & tō beranne \\
sęllan & tō sęllanne \\
\bottomrule
\end{tabular}
\par\endgroup

% mdtable:T26:end
% mdparagraph:P143:begin
The plain infinitive has no person or number ending. After \textbf{tō}, the selected form is inflected in \textbf{-anne}. Keep the doubled \textbf{ll} of \textbf{sęllanne}; the singular command’s consonant simplification is not an infinitive rule. [4]
% mdparagraph:P143:end

\end{minipage}\par\addvspace{4pt}

\par\begin{minipage}{\linewidth}\raggedright
% item: NEW-U06-N02
% key: KEY-NEW-U06-N02
% mdparagraph:P144:begin
\textbf{N02.}
% mdparagraph:P144:end

% mdparagraph:P145:begin
\textbf{Target: the plain and selected inflected infinitives.}
% mdparagraph:P145:end

% mdtable:T27:begin
\par\begingroup
\fontsize{10.5}{13.2}\selectfont
\setlength{\LTleft}{0pt}\setlength{\LTright}{\fill}
\noindent
\begin{tabular}{@{}>{\raggedright\arraybackslash}p{61.44mm}>{\raggedright\arraybackslash}p{59.75mm}@{}}
\toprule
\textbf{Plain infinitive} & \textbf{After tō} \\
\midrule
hīeran & tō hīeranne \\
cyssan & tō cyssanne \\
nęrian & tō nęrianne \\
gefręmman & tō gefręmmanne \\
gierwan & tō gierwanne \\
\bottomrule
\end{tabular}
\par\endgroup

% mdtable:T27:end
% mdparagraph:P146:begin
Retain the infinitive stems, including \textbf{ss, mm}, the \textbf{i} of \textbf{nęrianne}, and the \textbf{w} of \textbf{gierwanne}. The shorter singular imperative stems are not the basis for these forms. The task selects \textbf{-anne}; a documented \textbf{-enne} variant is valid language but differs from this explicit selection. [3–4]
% mdparagraph:P146:end

\end{minipage}\par\addvspace{4pt}

\par\begin{minipage}{\linewidth}\raggedright
% item: NEW-U06-N03
% key: KEY-NEW-U06-N03
% mdparagraph:P147:begin
\textbf{N03.}
% mdparagraph:P147:end

% mdparagraph:P148:begin
\textbf{Target: the plain and selected inflected infinitives.}
% mdparagraph:P148:end

% mdtable:T28:begin
\par\begingroup
\fontsize{10.5}{13.2}\selectfont
\setlength{\LTleft}{0pt}\setlength{\LTright}{\fill}
\noindent
\begin{tabular}{@{}>{\raggedright\arraybackslash}p{67.8mm}>{\raggedright\arraybackslash}p{53.39mm}@{}}
\toprule
\textbf{Plain infinitive} & \textbf{After tō} \\
\midrule
lufian & tō lufianne \\
habban & tō habbanne \\
libban & tō libbanne \\
sęcgan & tō sęcganne \\
hycgan & tō hycganne \\
\bottomrule
\end{tabular}
\par\endgroup

% mdtable:T28:end
% mdparagraph:P149:begin
Weak-II \textbf{lufian} gives \textbf{lufianne}, retaining \textbf{i}. The four weak-III forms retain their checked infinitive stems, including \textbf{bb} and \textbf{cg}. Do not form these infinitives from the exceptional singular commands. [3–4]
% mdparagraph:P149:end

\end{minipage}\par\addvspace{4pt}

\par\begin{minipage}{\linewidth}\raggedright
% item: NEW-U06-N04
% key: KEY-NEW-U06-N04
% mdparagraph:P150:begin
\textbf{N04.}
% mdparagraph:P150:end

% mdparagraph:P151:begin
\textbf{Target: the plain and selected inflected infinitives.}
% mdparagraph:P151:end

% mdtable:T29:begin
\par\begingroup
\fontsize{10.5}{13.2}\selectfont
\setlength{\LTleft}{0pt}\setlength{\LTright}{\fill}
\noindent
\begin{tabular}{@{}>{\raggedright\arraybackslash}p{67.8mm}>{\raggedright\arraybackslash}p{53.39mm}@{}}
\toprule
\textbf{Plain infinitive} & \textbf{After tō} \\
\midrule
wrēon & tō wrēonne \\
tēon & tō tēonne \\
sēon & tō sēonne \\
slēan & tō slēanne \\
fōn & tō fōnne \\
\bottomrule
\end{tabular}
\par\endgroup

% mdtable:T29:end
% mdparagraph:P152:begin
These forms retain the infinitive vowel and add the selected inflected ending. They do not restore the singular imperative’s \textbf{h}. In particular, long \textbf{ēo/\allowbreak{}ēa/\allowbreak{}ō} remains in \textbf{tēonne, sēonne, slēanne, fōnne}; the short vowels of \textbf{seoh/\allowbreak{}sleah} do not replace the infinitive vowels. [3–4]
% mdparagraph:P152:end

\end{minipage}\par\addvspace{4pt}

\par\begin{minipage}{\linewidth}\raggedright
% item: NEW-U06-N05
% key: KEY-NEW-U06-N05
% mdparagraph:P153:begin
\textbf{N05.}
% mdparagraph:P153:end

% mdparagraph:P154:begin
\textbf{Target: the plain and selected inflected infinitives.}
% mdparagraph:P154:end

% mdtable:T30:begin
\par\begingroup
\fontsize{10.5}{13.2}\selectfont
\setlength{\LTleft}{0pt}\setlength{\LTright}{\fill}
\noindent
\begin{tabular}{@{}>{\raggedright\arraybackslash}p{67.8mm}>{\raggedright\arraybackslash}p{53.39mm}@{}}
\toprule
\textbf{Plain infinitive} & \textbf{After tō} \\
\midrule
gān & tō gānne \\
dōn & tō dōnne \\
wesan & tō wesanne \\
bēon & tō bēonne \\
\bottomrule
\end{tabular}
\par\endgroup

% mdtable:T30:end
% mdparagraph:P155:begin
The anomalous infinitives retain their own stems. Keep long \textbf{ā, ō, ēo} in \textbf{gānne, dōnne, bēonne}, and the separate token \textbf{tō}. The two be-infinitives are not interchangeable labels for every finite or wish form with a similar spelling. [1,4]
% mdparagraph:P155:end

\end{minipage}\par\addvspace{4pt}

\par\begin{minipage}{\linewidth}\raggedright
% item: NEW-U06-N06
% key: KEY-NEW-U06-N06
% mdparagraph:P156:begin
\textbf{N06.}
% mdparagraph:P156:end

% mdparagraph:P157:begin
\textbf{Hē wille etan. Hē cōm tō etanne.}
% mdparagraph:P157:end

% mdparagraph:P158:begin
The first says “He wants to eat”: present \textbf{wille} takes plain \textbf{etan}. The second says “He came to eat”: preterite \textbf{cōm} is followed by the selected purpose phrase \textbf{tō etanne}, and the arriving man is also the eater. The purpose does not itself assert that he ate. A plain-purpose construction can be valid in Old English, but this question expressly requests the taught inflected purpose construction. [4]
% mdparagraph:P158:end

\end{minipage}\par\addvspace{4pt}

\par\begin{minipage}{\linewidth}\raggedright
% item: NEW-U06-N07
% key: KEY-NEW-U06-N07
% mdparagraph:P159:begin
\textbf{N07.}
% mdparagraph:P159:end

% mdparagraph:P160:begin
\textbf{Ðā gefēran cōmon tō helpanne ðǣm męnn.}
% mdparagraph:P160:end

% mdparagraph:P161:begin
The companions came for the purpose of helping the man. They are the subject of \textbf{cōmon} and the understood actors of \textbf{helpanne}. \textbf{Ðǣm męnn} is dative: it is the person helped, not the subject of the infinitive. Other defensible positions of that dative phrase may preserve the same relation. A bare-purpose infinitive can be valid language, but would miss this task’s explicit inflected-infinitive target. An accusative helped person would be a government error. [4]
% mdparagraph:P161:end

\end{minipage}\par\addvspace{4pt}

\par\begin{minipage}{\linewidth}\raggedright
% item: NEW-U06-N08
% key: KEY-NEW-U06-N08
% mdparagraph:P162:begin
\textbf{N08.}
% mdparagraph:P162:end

% mdparagraph:P163:begin
\textbf{Se mann cōm tō giefanne męte ðǣm fuglum.}
% mdparagraph:P163:end

% mdparagraph:P164:begin
The man came in order to give food to the birds. The same man is both the arriving person and the intended giver. \textbf{Męte} is accusative; \textbf{ðǣm fuglum} is dative. The purpose construction uses \textbf{tō giefanne}. The separately taught giving verb \textbf{sęllan}, with its correct inflected infinitive, can preserve the same meaning. A defensible order of gift and recipient is acceptable; reversing their grammatical roles is not. [4]
% mdparagraph:P164:end

\end{minipage}\par\addvspace{4pt}

\par\begin{minipage}{\linewidth}\raggedright
% item: NEW-U06-N09
% key: KEY-NEW-U06-N09
% mdparagraph:P165:begin
\textbf{N09.}
% mdparagraph:P165:end

% mdtable:T31:begin
\par\begingroup
\fontsize{10.5}{13.2}\selectfont
\setlength{\LTleft}{0pt}\setlength{\LTright}{\fill}
\noindent
\begin{tabular}{@{}>{\raggedright\arraybackslash}p{26.82mm}>{\raggedright\arraybackslash}p{37.65mm}>{\raggedright\arraybackslash}p{25.74mm}>{\raggedright\arraybackslash}p{25.36mm}@{}}
\toprule
\textbf{Sentence} & \textbf{Finite verb} & \textbf{Infinitive} & \textbf{Construction} \\
\midrule
Hīe sculon ðæt weorc dōn. & sculon: present indicative, third plural & dōn: plain & Obligation with sculan \\
Hīe cōmon tō dōnne ðæt weorc. & cōmon: preterite indicative, third plural & dōnne: inflected after tō & Purpose of coming \\
\bottomrule
\end{tabular}
\par\endgroup

% mdtable:T31:end
% mdparagraph:P166:begin
“They must do the work.” “They came to do the work.” In both sentences the work is accusative and the subject is also the intended doer. The first imposes an obligation; the second states a past arrival and its purpose. Neither infinitive is another finite past form, and neither clause by itself asserts completion of the work. “Were obliged” would incorrectly move the first finite verb into the past. [4–5]
% mdparagraph:P166:end

\end{minipage}\par\addvspace{4pt}

\par\begin{minipage}{\linewidth}\raggedright
% item: NEW-U06-N10
% key: KEY-NEW-U06-N10
% mdparagraph:P167:begin
\textbf{N10.}
% mdparagraph:P167:end

% mdparagraph:P168:begin
\textbf{Wē magon ðone męte beran. Ic mæg ðone męte beran.}
% mdparagraph:P168:end

% mdparagraph:P169:begin
The meanings are “We can carry the food” and “I can carry the food.” The finite modal changes from first-plural \textbf{magon} to first-singular \textbf{mæg}. Plain \textbf{beran} remains the same, as does its accusative food object. Changing the infinitive to an agreeing finite form, or changing the object merely because the subject becomes singular, would miss the construction.
% mdparagraph:P169:end

\end{minipage}\par\addvspace{4pt}

\reserve{7\baselineskip}
% mdparagraph:P170:begin
\section*{C. Food and water for the journey}
% mdparagraph:P170:end

\par\begin{minipage}{\linewidth}\raggedright
% item: NEW-U06-C01
% key: KEY-NEW-U06-C01
% mdparagraph:P171:begin
\textbf{C01.}
% mdparagraph:P171:end

% mdparagraph:P172:begin
The man asks for help. \textbf{Two companions} come, and their stated purpose is \textbf{to help}.
% mdparagraph:P172:end

% mdparagraph:P173:begin
The number comes from \textbf{twēgen gefēran}; the purpose comes from \textbf{tō helpanne}. It is reasonable to understand that they are helping the man who called them, but the purpose phrase itself has no expressed object. “They came and helped” can suggest a completed action that the purpose phrase alone does not state. The later narrative supplies the actions they actually perform.
% mdparagraph:P173:end

\end{minipage}\par\addvspace{4pt}

\par\begin{minipage}{\linewidth}\raggedright
% item: NEW-U06-C02
% key: KEY-NEW-U06-C02
% mdparagraph:P174:begin
\textbf{C02.}
% mdparagraph:P174:end

% mdparagraph:P175:begin
\textbf{Target: all seven commands, their plural addressees, and the corresponding singulars.}
% mdparagraph:P175:end

% mdtable:T32:begin
\par\begingroup
\fontsize{10.5}{13.2}\selectfont
\setlength{\LTleft}{0pt}\setlength{\LTright}{\fill}
\noindent
\begin{tabular}{@{}>{\raggedright\arraybackslash}p{40.6mm}>{\raggedright\arraybackslash}p{39.79mm}>{\raggedright\arraybackslash}p{37.99mm}@{}}
\toprule
\textbf{Command in the text} & \textbf{Addressee number} & \textbf{Singular command} \\
\midrule
Cumað & Plural & Cum \\
helpað & Plural & help \\
Nimað & Plural & Nim \\
berað & Plural & ber \\
Giefað & Plural & Gief \\
Sittað & Plural & Site \\
etað & Plural & et \\
\bottomrule
\end{tabular}
\par\endgroup

% mdtable:T32:end
% mdparagraph:P176:begin
All address the companions, so all are second-person plural. The list does not include finite statements such as \textbf{wille, mæg} or \textbf{sculon}, nor the purpose infinitive \textbf{helpanne}. \textbf{Site} needs its final \textbf{-e} and single \textbf{t}; it is not made by merely shortening the plural to its doubled stem.
% mdparagraph:P176:end

\end{minipage}\par\addvspace{4pt}

\par\begin{minipage}{\linewidth}\raggedright
% item: NEW-U06-C03
% key: KEY-NEW-U06-C03
% mdparagraph:P177:begin
\textbf{C03.}
% mdparagraph:P177:end

% mdparagraph:P178:begin
\textbf{Hine} is masculine accusative singular and most naturally refers to \textbf{ðone męte}, “the food,” which has just been named. English “it” does not make the Old English pronoun neuter. Context, as well as grammatical compatibility, identifies the food; the form alone does not exclude every other masculine referent in the passage.
% mdparagraph:P178:end

% mdparagraph:P179:begin
The \textbf{birds} explicitly receive the water: they are the dative recipients \textbf{ðǣm fuglum}. The companions subsequently \textbf{eat}. It is natural to infer that they eat the mentioned food, but \textbf{ǣton} has no expressed object, so “they ate that food” is an inference rather than an additional explicit clause. The final \textbf{hīe} likewise need not uniquely encode whether only the companions or the whole party travelled. Answers should preserve these differences between what is stated and what context suggests.
% mdparagraph:P179:end

\end{minipage}\par\addvspace{4pt}

\par\begin{minipage}{\linewidth}\raggedright
% item: NEW-U06-C04
% key: KEY-NEW-U06-C04
% mdparagraph:P180:begin
\textbf{C04.}
% mdparagraph:P180:end

% mdtable:T33:begin
\par\begingroup
\fontsize{10.5}{13.2}\selectfont
\setlength{\LTleft}{0pt}\setlength{\LTright}{\fill}
\noindent
\begin{tabular}{@{}>{\raggedright\arraybackslash}p{17.41mm}>{\raggedright\arraybackslash}p{24.9mm}>{\raggedright\arraybackslash}p{23.68mm}>{\raggedright\arraybackslash}p{24.64mm}>{\raggedright\arraybackslash}p{22.12mm}@{}}
\toprule
\textbf{Form} & \textbf{Person and number here} & \textbf{Tense and mood} & \textbf{Stated sense} & \textbf{Plain infinitive} \\
\midrule
wille & First singular & Present indicative & Want/\allowbreak{}intend to & faran \\
mæg & First singular & Present indicative & Be able to & beran \\
sculon & First plural & Present indicative & Must/\allowbreak{}be obliged to & gefręmman \\
\bottomrule
\end{tabular}
\par\endgroup

% mdtable:T33:end
% mdparagraph:P181:begin
The \textbf{ic} before \textbf{wille} belongs to the man speaking; the \textbf{ic} before \textbf{mæg} belongs to the companion speaking. Neither is third singular simply because an English translation elsewhere uses “he can” or “he wants.” \textbf{Wē} selects first plural for \textbf{sculon}. The following infinitives name the actions and carry no personal agreement.
% mdparagraph:P181:end

\end{minipage}\par\addvspace{4pt}

\par\begin{minipage}{\linewidth}\raggedright
% item: NEW-U06-C05
% key: KEY-NEW-U06-C05
% mdparagraph:P182:begin
\textbf{C05.}
% mdparagraph:P182:end

% mdparagraph:P183:begin
The man said, “I want to travel. My ship is ready. Come and help me.” Two companions came to help. The man said, “Take the food and carry it. Give the birds the water.” The companion said, “I can carry the food. We must perform the work.” The companions carried the food and gave the birds the water. The man said, “Sit and eat.” The companions ate. They travelled.
% mdparagraph:P183:end

% mdparagraph:P184:begin
\textbf{Essential meaning:} retain the man’s wish, the companion’s ability, the group’s obligation, the plural commands, the purpose of coming, and the narrated past carrying, giving, eating and travelling. “Go” for \textbf{faran}, “do/\allowbreak{}accomplish the work” for \textbf{gefręmman}, and “are obliged to” for \textbf{sculon} are defensible. The quoted \textbf{sculon} is present: do not shift it into the past merely because the quotation is introduced by \textbf{cwæð}.
% mdparagraph:P184:end

% mdparagraph:P185:begin
Food is the natural referent of \textbf{hine}, and eating that food is inferred, as C03 explains. Either a companions-only or a contextually inclusive reading of final “they” can be defended. The text says that the ship is ready and that they travel; it does not explicitly narrate boarding, sailing, or a destination. Translating willingness as an accomplished departure, or the purpose of helping as proof that helping has already occurred, would add to the stated meaning.
% mdparagraph:P185:end

\end{minipage}\par\addvspace{4pt}

\par\begin{minipage}{\linewidth}\raggedright
% item: NEW-U06-C06
% key: KEY-NEW-U06-C06
% mdparagraph:P186:begin
\textbf{C06.}
% mdparagraph:P186:end

% mdparagraph:P187:begin
\textbf{Commands to one companion:} \textbf{Cum and help mē.}
% mdparagraph:P187:end

% mdparagraph:P188:begin
\textbf{First narrative sentence in the present, with its quotation unchanged:}
% mdparagraph:P188:end

% mdparagraph:P189:begin
\textbf{Se mann cwið: ‘Ic wille faran. Mīn scip is gearu. Cumað and helpað mē.’}
% mdparagraph:P189:end

% mdparagraph:P190:begin
The first operation changes the command addressee from plural to singular. The second changes only the reporting \textbf{cwæð}, “said,” to present \textbf{cwið}, “says.” It is a separate operation: the quoted commands stay plural in that second answer, as the instruction requires. Do not carry the first answer’s singular commands into the quotation that was expressly to remain unchanged.
% mdparagraph:P190:end

\end{minipage}\par\addvspace{4pt}

\reserve{7\baselineskip}
% mdparagraph:P191:begin
\section*{R. The fisherman’s gift and first direction}
% mdparagraph:P191:end

\par\begin{minipage}{\linewidth}\raggedright
% item: NEW-U06-R01
% key: KEY-NEW-U06-R01
% mdparagraph:P192:begin
\textbf{R01.}
% mdparagraph:P192:end

% mdparagraph:P193:begin
The fisherman sees the young man lying at his feet and lifts him up with compassion. He divides \textbf{his cloak} into two and gives Apollonius \textbf{one half}. He has also set food before him earlier in the excerpt.
% mdparagraph:P193:end

% mdparagraph:P194:begin
The target is the supported sequence of care and gift. A natural answer may call the young man Apollonius, using the context. It must not say that the fisherman divides the food or gives away the entire cloak. The selected text explicitly makes the sight of the young man’s posture the immediate occasion for the action. [8]
% mdparagraph:P194:end

\end{minipage}\par\addvspace{4pt}

\par\begin{minipage}{\linewidth}\raggedright
% item: NEW-U06-R02
% key: KEY-NEW-U06-R02
% mdparagraph:P195:begin
\textbf{R02.}
% mdparagraph:P195:end

% mdparagraph:P196:begin
\textbf{Wolde} is preterite indicative third singular of \textbf{willan}, with the fisherman as its subject. Its plain infinitive is the supplied \textbf{gecȳðan}, “show.” He wanted to show Apollonius greater kindness, according to his ability or means.
% mdparagraph:P196:end

% mdparagraph:P197:begin
“Wished” or “intended” can preserve the volitional meaning; an automatic hypothetical “would have shown” would add a condition. The isolated spelling \textbf{wolde} could also serve first singular, but \textbf{hē} fixes third singular here. Supplied \textbf{mihte} in \textbf{be his mihte} is the noun “ability/\allowbreak{}power,” not another finite modal verb. [8]
% mdparagraph:P197:end

\end{minipage}\par\addvspace{4pt}

\par\begin{minipage}{\linewidth}\raggedright
% item: NEW-U06-R03
% key: KEY-NEW-U06-R03
% mdparagraph:P198:begin
\textbf{R03.}
% mdparagraph:P198:end

% mdtable:T34:begin
\par\begingroup
\fontsize{10.5}{13.2}\selectfont
\setlength{\LTleft}{0pt}\setlength{\LTright}{\fill}
\noindent
\begin{tabular}{@{}>{\raggedright\arraybackslash}p{28.55mm}>{\raggedright\arraybackslash}p{35.37mm}>{\raggedright\arraybackslash}p{54.45mm}@{}}
\toprule
\textbf{Source phrase} & \textbf{Plain infinitive} & \textbf{Form and local meaning} \\
\midrule
tō bēodanne & bēodan & Inflected infinitive: to offer \\
tō sęllanne & sęllan & Inflected infinitive: to give \\
\bottomrule
\end{tabular}
\par\endgroup

% mdtable:T34:end
% mdparagraph:P199:begin
The fisherman sets out the food \textbf{he had available to offer} Apollonius. In his speech, he says to take \textbf{what he has available to give} him. The persons who receive the offering/\allowbreak{}gift are supplied by \textbf{him} and \textbf{ðē} in their respective clauses.
% mdparagraph:P199:end

% mdparagraph:P200:begin
The target forms are inflected infinitives after \textbf{tō}, not finite verbs with person or tense. Use the printed local explanation for the \textbf{habban} construction. Modern English “had to/\allowbreak{}must” would introduce an obligation that these availability phrases do not express. The supplied relative clauses need not be independently reconstructed. [4,8]
% mdparagraph:P200:end

\end{minipage}\par\addvspace{4pt}

\par\begin{minipage}{\linewidth}\raggedright
% item: NEW-U06-R04
% key: KEY-NEW-U06-R04
% mdparagraph:P201:begin
\textbf{R04.}
% mdparagraph:P201:end

% mdtable:T35:begin
\par\begingroup
\fontsize{10.5}{13.2}\selectfont
\setlength{\LTleft}{0pt}\setlength{\LTright}{\fill}
\noindent
\begin{tabular}{@{}>{\raggedright\arraybackslash}p{24.19mm}>{\raggedright\arraybackslash}p{23.83mm}>{\raggedright\arraybackslash}p{43.37mm}>{\raggedright\arraybackslash}p{24.19mm}@{}}
\toprule
\textbf{Source command} & \textbf{Verb} & \textbf{Analysis in this speech} & \textbf{Plural command} \\
\midrule
Nim & niman & Imperative, second singular & Nimað \\
gā & gān & Imperative, second singular & gāð \\
\bottomrule
\end{tabular}
\par\endgroup

% mdtable:T35:end
% mdparagraph:P202:begin
The fisherman directly addresses Apollonius: “Take … and go ….” That speech context establishes command use. \textbf{Gā} is also the spelling of a first-singular indicative cell elsewhere, but it is not an assertion “I go” here. Both plural replacements would address more than one person. [1,8]
% mdparagraph:P202:end

\end{minipage}\par\addvspace{4pt}

\par\begin{minipage}{\linewidth}\raggedright
% item: NEW-U06-R05
% key: KEY-NEW-U06-R05
% mdparagraph:P203:begin
\textbf{R05.}
% mdparagraph:P203:end

% mdparagraph:P204:begin
\textbf{Cweðende}, “saying,” is the supplied present active participle accompanying the fisherman’s action. It is the form that introduces his actual words, with \textbf{ðus}, “thus.” It is not a finite past or a command.
% mdparagraph:P204:end

% mdparagraph:P205:begin
\textbf{Cwæð} is a finite preterite indicative, “said,” first or third singular when isolated; \textbf{gā} in this speech is a second-singular command, “go.” An English rendering such as “and said” may express the participial event naturally, but it does not make the Old English \textbf{cweðende} finite. The local note supplies the participial construction; no independent participle formation is required. [8]
% mdparagraph:P205:end

\end{minipage}\par\addvspace{4pt}

\par\begin{minipage}{\linewidth}\raggedright
% item: NEW-U06-R06
% key: KEY-NEW-U06-R06
% mdparagraph:P206:begin
\textbf{R06.}
% mdparagraph:P206:end

% mdparagraph:P207:begin
He says he has something available to give Apollonius, understood in context as the half-cloak just given. He tells Apollonius to \textbf{take it and go into the town}.
% mdparagraph:P207:end

% mdparagraph:P208:begin
The supported \textbf{ðæt ic ðē tō sęllanne hæbbe} means “what/\allowbreak{}that which I have to give you.” “What” introduces the thing available for giving; it is not a question asking what the fisherman possesses. \textbf{Hæbbe} supplies “I have” in that local explanation. Its spelling alone cannot uniquely distinguish every possible mood, so the answer should not manufacture a mood test or an additional command “have!” from it. [8]
% mdparagraph:P208:end

\end{minipage}\par\addvspace{4pt}

\par\begin{minipage}{\linewidth}\raggedright
% item: NEW-U06-R07
% key: KEY-NEW-U06-R07
% mdparagraph:P209:begin
\textbf{R07.}
% mdparagraph:P209:end

% mdtable:T36:begin
\par\begingroup
\fontsize{10.5}{13.2}\selectfont
\setlength{\LTleft}{0pt}\setlength{\LTright}{\fill}
\noindent
\begin{tabular}{@{}>{\raggedright\arraybackslash}p{31.95mm}>{\raggedright\arraybackslash}p{51.83mm}>{\raggedright\arraybackslash}p{34.6mm}@{}}
\toprule
\textbf{Phrase} & \textbf{Kind} & \textbf{Local meaning} \\
\midrule
tō bēodanne & tō with inflected infinitive & to offer \\
tō sęllanne & tō with inflected infinitive & to give \\
tō his hūse & Destination with a noun phrase & to his house \\
intō ðǣre ceastre & Destination with a noun phrase & into the town/\allowbreak{}city \\
\bottomrule
\end{tabular}
\par\endgroup

% mdtable:T36:end
% mdparagraph:P210:begin
The first two contain the inflected verb forms. \textbf{Hūse} and \textbf{ceastre} are supplied nouns, both dative singular in these actual phrases; \textbf{ðǣre} is feminine dative singular. A shared sequence of letters \textbf{tō} does not make every phrase infinitival. Preserve the local dative after \textbf{intō}: motion is not a reason to replace the printed form with an accusative or derive a universal rule from this one phrase. [8]
% mdparagraph:P210:end

\end{minipage}\par\addvspace{4pt}

\par\begin{minipage}{\linewidth}\raggedright
% item: NEW-U06-R08
% key: KEY-NEW-U06-R08
% mdparagraph:P211:begin
\textbf{R08.}
% mdparagraph:P211:end

% mdparagraph:P212:begin
Then, as soon as the fisherman saw that the young man was lying at his feet, he compassionately lifted him up and led him with him to his house, and set before him the provisions that he had available to offer him. He still wanted, according to his ability, to show him greater kindness; he then tore his cloak in two and gave Apollonius one half, saying: “Take what I have to give you, and go into the town.”
% mdparagraph:P212:end

% mdparagraph:P213:begin
\textbf{Essential meaning:} preserve the opening temporal and content relations, the fisherman’s compassionate action, the leading home and provision of food, his wish to show greater kindness, the division and gift of the cloak, and both final commands. The local help remains available throughout.
% mdparagraph:P213:end

% mdparagraph:P214:begin
“With compassion” is a closer rendering of \textbf{mid mildheortnesse}; “compassionately” is natural English. “Food/\allowbreak{}provisions,” “still/\allowbreak{}further,” “kindness/\allowbreak{}compassion,” “one half/\allowbreak{}the half part,” and “town/\allowbreak{}city” are defensible within the stated relations. “As far as he could” can render \textbf{be his mihte} naturally, provided it is understood as the supplied phrase rather than a second finite modal in the Old English. Context identifies the fisherman’s feet/\allowbreak{}house and Apollonius as the beneficiary; the pronoun spellings alone do not provide all those identities.
% mdparagraph:P214:end

% mdparagraph:P215:begin
Do not translate the availability phrases as obligations, render \textbf{cweðende} as a command, or turn the final \textbf{gā} into a past departure. Stop after the first direction sentence. This excerpt does not narrate its fulfillment or contain the continuation of the fisherman’s speech. [8]
% mdparagraph:P215:end

\end{minipage}\par\addvspace{4pt}

\reserve{7\baselineskip}
% mdparagraph:P216:begin
\section*{X. Cumulative practice}
% mdparagraph:P216:end

\par\begin{minipage}{\linewidth}\raggedright
% item: NEW-U06-X01
% key: KEY-NEW-U06-X01
% mdparagraph:P217:begin
\textbf{X01.}
% mdparagraph:P217:end

% mdparagraph:P218:begin
\textbf{Wit woldon uncerne męte ðǣm fuglum giefan. Wit gēafon hine him.}
% mdparagraph:P218:end

% mdparagraph:P219:begin
The first sentence states the two speakers’ past wish to give their food to the birds; the second states that the two speakers actually gave it. \textbf{Wit} preserves the dual subject. \textbf{Uncerne} is masculine accusative singular, agreeing with \textbf{męte}, while its possessive stem identifies the two possessors. \textbf{Hine} refers to that masculine food; \textbf{him} is dative plural here, referring to the birds.
% mdparagraph:P219:end

% mdparagraph:P220:begin
The dual subjects take the plural finite forms \textbf{woldon/\allowbreak{}gēafon}. The infinitive remains plain after \textbf{woldon}. A defensible order or another taught giving verb may preserve the meaning, but ordinary plural possessors would lose the requested dual relation, and a neuter “it” would lose the food’s grammatical gender.
% mdparagraph:P220:end

\end{minipage}\par\addvspace{4pt}

\par\begin{minipage}{\linewidth}\raggedright
% item: NEW-U06-X02
% key: KEY-NEW-U06-X02
% mdparagraph:P221:begin
\textbf{X02.}
% mdparagraph:P221:end

% mdparagraph:P222:begin
\textbf{Se ealda mann ne mæg ðæt micle scip beran. Helpað him.}
% mdparagraph:P222:end

% mdparagraph:P223:begin
The man lacks the ability to carry the large ship; several companions are told to help him. \textbf{Ealda} is weak nominative masculine singular; \textbf{micle} is weak accusative neuter singular. Each adjective belongs to its own article-and-noun phrase. \textbf{Him} is dative singular as the person helped, while \textbf{helpað} is plural because of its addressees.
% mdparagraph:P223:end

% mdparagraph:P224:begin
A defensible word order can vary. The man’s singular number does not make the command singular: he is its helped person, not its addressee. Nor does \textbf{helpan} take an accusative merely because \textbf{beran} does.
% mdparagraph:P224:end

\end{minipage}\par\addvspace{4pt}

\par\begin{minipage}{\linewidth}\raggedright
% item: NEW-U06-X03
% key: KEY-NEW-U06-X03
% mdparagraph:P225:begin
\textbf{X03.}
% mdparagraph:P225:end

% mdparagraph:P226:begin
\textbf{Hē cōm tō helpanne ūs. Wē woldon him męte giefan. Hē sæt and ǣt.}
% mdparagraph:P226:end

% mdparagraph:P227:begin
He came for the purpose of helping us; we wanted to give him food; he sat and ate. The arriving man is the understood helper. \textbf{Ūs} is dative after \textbf{helpan}; \textbf{him} is the dative gift-recipient and \textbf{męte} the accusative gift. The first infinitive is inflected after \textbf{tō}, the second plain after \textbf{woldon}.
% mdparagraph:P227:end

% mdparagraph:P228:begin
The final coordinated verbs are both third-singular preterites. Preserve short \textbf{æ} in \textbf{sæt} and long \textbf{ǣ} in \textbf{ǣt}. Other defensible word orders or the taught alternative giving verb may preserve the sense, but the purpose must not be mistaken for a completed helping event, and the wish to give must not be mistaken for the giving itself. [4]
% mdparagraph:P228:end

\end{minipage}\par\addvspace{4pt}

\par\begin{minipage}{\linewidth}\raggedright
% item: NEW-U06-X04
% key: KEY-NEW-U06-X04
% mdparagraph:P229:begin
\textbf{X04.}
% mdparagraph:P229:end

% mdparagraph:P230:begin
\textbf{Gē gāð} is a statement: “You, plural, go.” \textbf{Gāð} is a command: “Go,” addressed to more than one person. In the first, \textbf{gē} fixes second-plural agreement; the task’s stipulated statement use identifies \textbf{gāð} as present indicative. An expressed second-person subject can also accompany a command. The second uses the second-plural imperative. Their spelling is shared, but their functions differ.
% mdparagraph:P230:end

% mdparagraph:P231:begin
\textbf{Hē dēð ðæt weorc} means “He does the work,” with third-singular present indicative \textbf{dēð}. \textbf{Dō ðæt weorc} means “Do the work,” a command to one person. Here the indicative and imperative forms visibly differ. \textbf{Dō} can serve first-singular indicative in a different sentence, so an ending alone is not a complete account of its use. Identify the actual statement or command in context. [1,3]
% mdparagraph:P231:end

\end{minipage}\par\addvspace{4pt}

\par\begin{minipage}{\linewidth}\raggedright
% item: NEW-U06-X05
% key: KEY-NEW-U06-X05
% mdparagraph:P232:begin
\textbf{X05.}
% mdparagraph:P232:end

% mdparagraph:P233:begin
\textbf{One model:}
% mdparagraph:P233:end

% mdparagraph:P234:begin
\textbf{Se mann cōm tō helpanne twǣm gefērum. Ðā gefēran willað ðone męte beran. Berað ðone męte. Hīe bǣron ðone męte.}
% mdparagraph:P234:end

% mdparagraph:P235:begin
“The man came to help two companions. The companions want to carry the food. Carry the food. They carried the food.”
% mdparagraph:P235:end

% mdtable:T37:begin
\par\begingroup
\fontsize{10.5}{13.2}\selectfont
\setlength{\LTleft}{0pt}\setlength{\LTright}{\fill}
\noindent
\begin{tabular}{@{}>{\raggedright\arraybackslash}p{40.74mm}>{\raggedright\arraybackslash}p{80.45mm}@{}}
\toprule
\textbf{Required part} & \textbf{How the model meets it} \\
\midrule
Past arrival with purpose & cōm tō helpanne: the arriving man intends to help \\
Willingness & willað … beran: a finite volitional verb with plain infinitive \\
Command to both companions & Berað: second-person plural imperative \\
Completed past action & bǣron: the companions carried the food \\
\bottomrule
\end{tabular}
\par\endgroup

% mdtable:T37:end
% mdparagraph:P236:begin
The three new constructions to label are \textbf{purpose with tō + inflected infinitive}, \textbf{willan + plain infinitive}, and \textbf{the plural direct command}. \textbf{Twǣm gefērum} is dative because the companions are helped; \textbf{ðone męte} is accusative because the food is carried.
% mdparagraph:P236:end

% mdparagraph:P237:begin
Give credit to another coherent four-sentence account that uses the requested vocabulary and constructions, makes the two-companion situation clear, and preserves the stated distinction between intention and accomplished action. The arriving/\allowbreak{}helper roles or the order of phrases can differ from this model when the reference remains coherent. Do not require this exact string or an extra narrative event. In particular, check that the arriving subject is also the intended actor of the purpose infinitive, that the volitional complement is plain, that the command addresses both companions with a plural form, and that the last required event has an actual finite past verb. Correct vocabulary alone does not repair a wrong case, command number, or infinitive construction.
% mdparagraph:P237:end

\end{minipage}\par\addvspace{4pt}

\clearpage
% mdparagraph:P238:begin
\section*{Sources}
% mdparagraph:P238:end

\begingroup\fontsize{10}{12.8}\selectfont
% mdparagraph:P239:begin
PDF page numbers below count from the first PDF page, including the preliminary pages. Editorial line numbers count the printed lines of a selected page; they are not additional numbers printed by Cook.
% mdparagraph:P239:end

% mdparagraph:P240:begin
[1] Albert S. Cook, \textbf{A First Book in Old English}, second edition (1895), §§125–137, pp. 76–78 (PDF pp. 92–94), for the preterite-presents; §§138–141, pp. 78–80 (PDF pp. 94–96), for \textbf{bēon, dōn, gān} and \textbf{willan}. The tables above retain the unit’s stated choices among the printed variants.
% mdparagraph:P240:end

% mdparagraph:P241:begin
[2] R. D. Fulk, \textbf{An Introductory Grammar of Old English with an Anthology of Readings} (2020), §52, p. 33 (PDF p. 45), on the temporal and habitual uses of the be-series; §123, pp. 89–90 (PDF pp. 101–102), on the preterite-present systems and the limits of citation forms.
% mdparagraph:P241:end

% mdparagraph:P242:begin
[3] Cook, §§93–95, pp. 54–55 (PDF pp. 70–71), and §§100–101, pp. 57–59 (PDF pp. 73–75), for strong and contracted forms; §§106–107, pp. 64–65 (PDF pp. 80–81), for the j-presents; §§112–124, pp. 68–76 (PDF pp. 84–92), for the weak systems; §35, pp. 24–25 (PDF pp. 40–41), for final consonant simplification. Compare Joseph Wright and Elizabeth Mary Wright, \textbf{Old English Grammar}, second edition (1914), §§478, 480, pp. 250–251 (PDF pp. 270–271), and §§521–525, 529–530, pp. 266–268, 270–271 (PDF pp. 286–288, 290–291). Fulk, §127, pp. 93–94 (PDF pp. 105–106), and §134, p. 102 (PDF p. 114), corroborates the j-present and anomalous distinctions.
% mdparagraph:P242:end

% mdparagraph:P243:begin
[4] Cook, §§199.1–200, p. 106 (PDF p. 122), distinguishes the bare purpose infinitive from the inflected construction. Wright and Wright, §§478, 480, pp. 250–251 (PDF pp. 270–271), and Fulk, §58, p. 38 (PDF p. 50), support the infinitive forms and ending variants.
% mdparagraph:P243:end

% mdparagraph:P244:begin
[5] Joseph Bosworth and T. Northcote Toller, \textbf{An Anglo-Saxon Dictionary} (main dictionary, 1898): \textbf{CUNNAN}, p. 174b, II; \textbf{MAGAN}, pp. 664b–665b, III(1); \textbf{ÞURFAN}, p. 1076b, II; \textbf{DURRAN}, p. 219b; \textbf{SCULAN}, pp. 843b–844a, II; \textbf{MÓTAN}, p. 699a, I; \textbf{WILLAN}, p. 1226b, II(a); \textbf{DÓN}, p. 208b; \textbf{GÁN}, p. 360b; \textbf{FARAN}, p. 270b, I. These locators support the bounded senses used in the models, without limiting each verb to a single meaning everywhere.
% mdparagraph:P244:end

% mdparagraph:P245:begin
[6] Cook, Reader XII, p. 173 (PDF p. 189), printed line 5 and editorial lines 13–14; Bosworth–Toller, \textbf{WITAN}, p. 1244a, I(2), and p. 1244b, I(6)(a), for knowledge with an accusative nominal object and the fact/\allowbreak{}matter use of \textbf{ðing}.
% mdparagraph:P245:end

% mdparagraph:P246:begin
[7] Cook, §198, p. 105 (PDF p. 121), on direct commands. For \textbf{ne} with an imperative, compare Reader XII, p. 183 (PDF p. 199), editorial line 7, and the cited examples in §156h, p. 92 (PDF p. 108), and §159a, p. 93 (PDF p. 109).
% mdparagraph:P246:end

% Keep the complete final source note and its original-note quotations together.
\reserve{10\baselineskip}
% mdparagraph:P247:begin
[8] Cook, Reader XII, p. 167 (PDF p. 183), editorial lines 3–10, spanning the printed line marks 5 and 10. The selected excerpt ends at the period after \textbf{ceastre}, completing the fisherman’s first direction sentence; his speech continues beyond the selection. Cook’s original notes are 5, “See 143.”; 6, “Lat. epulas.”; and 7, “MS. fæstnesse. Lat. pietati.” The local support for the temporal clause, relative phrases, availability construction and participle remains part of the reading. The two \textbf{tō}-infinitives can be identified from the taught forms without turning that availability construction into an independently required rule.
% mdparagraph:P247:end

\endgroup
\end{document}
